\documentclass[10pt]{article}
\usepackage[a4paper,margin=0.82in]{geometry}
\usepackage{amsmath,amssymb,amsthm,mathtools,amscd,enumitem,xcolor,hyperref}
\hypersetup{colorlinks=true,linkcolor=blue!55!black,citecolor=blue!55!black,urlcolor=blue!55!black}
\setlength{\emergencystretch}{4em}
\setlength{\parindent}{0pt}
\setlength{\parskip}{0.5em}
\setlist{leftmargin=1.6em,itemsep=0.2em,topsep=0.25em}
\sloppy

\newtheorem{theorem}{Theorem}[section]
\newtheorem{proposition}[theorem]{Proposition}
\newtheorem{lemma}[theorem]{Lemma}
\newtheorem{corollary}[theorem]{Corollary}
\theoremstyle{definition}
\newtheorem{definition}[theorem]{Definition}
\newtheorem{remark}[theorem]{Remark}

\newcommand{\Ab}{\mathsf{Ab}}
\newcommand{\Mod}{\mathsf{Mod}}
\newcommand{\Sh}{\mathsf{Sh}}
\newcommand{\Spec}{\operatorname{Spec}}
\newcommand{\Hom}{\operatorname{Hom}}
\newcommand{\colim}{\operatorname*{colim}}
\newcommand{\Rlim}{R\!\varprojlim}
\newcommand{\RGamma}{R\Gamma}
\newcommand{\F}{\mathbf F}
\newcommand{\Zl}{\mathbf Z_\ell}
\newcommand{\Ql}{\mathbf Q_\ell}
\newcommand{\Qlb}{\overline{\mathbf Q}_\ell}
\newcommand{\Dbc}{D^b_c}
\newcommand{\Dplus}{D^+}
\newcommand{\Et}{\mathrm{et}}
\newcommand{\Proet}{\mathrm{proet}}
\newcommand{\Ex}{\operatorname{Ex}}
\newcommand{\deps}[1]{\par\smallskip\textit{Direct dependencies.} #1\par}
\newcommand{\sources}[1]{\par\textit{Source locator.} #1\par}

\pdfstringdefDisableCommands{%
  \def\Qlb{Qbar_l}%
}

\begin{document}
\begin{center}
{\large\bfseries Proper Base Change\par}
\medskip
{\normalsize From finite morphisms to complete pro-etale coefficients\par}
\end{center}

\begin{abstract}
The proper-base-change theorem for torsion coefficients is assembled from
finite morphisms, composition and cancellation, Chow reduction, the
projective line over a strictly henselian base, and invariance for proper
curves under separably closed extension.  It then extends to bounded-below
complexes and passes from the etale site to the pro-etale site.  Derived
completion and finite reduction give the corresponding theorem for complete
adic coefficients, finite extensions of $\Ql$, and the Karoubi-completed
$\Qlb$ category.  All derived categories below are ordinary triangulated
categories.
\end{abstract}

\tableofcontents

\section{Ambient categories and the exchange morphism}

\begin{definition}[Sheaf and derived categories]
For a scheme $T$, let $T_{\Et}$ be the small etale site and let
$\Ab(T_{\Et})$ be the Grothendieck abelian category of sheaves of abelian
groups.  For a commutative ring $R$, let $R_T$ be the constant etale sheaf
of rings and write $\Mod(T_{\Et},R)$ for sheaves of $R_T$-modules.  Define
the ordinary Verdier derived categories
\[
 D(T_{\Et},R)=
 \mathsf{Cpx}(\Mod(T_{\Et},R))
 [\{\text{quasi-isomorphisms}\}^{-1}],
 \qquad \Dplus(T_{\Et},R)\subset D(T_{\Et},R).
\]
The second category consists of objects whose cohomology sheaves vanish
below some integer.  The notation without $R$ refers to sheaves of abelian
groups.  An object has torsion cohomology if every local cohomology section
is killed by a positive integer.
\end{definition}
\deps{The small etale site, sheaves of modules, cochain complexes, and
Verdier localization.}

\begin{definition}[Direct and inverse image]
For a morphism $u:U\to V$, direct image is the left exact functor
\[
 u_*:\Ab(U_{\Et})\longrightarrow\Ab(V_{\Et}),
 \qquad (u_*F)(W)=F(W\times_VU).
\]
If $K\to I(K)$ is a $K$-injective resolution, set
\[
 Ru_*K=u_*I(K),\qquad
 Ru_*:D(U_{\Et})\longrightarrow D(V_{\Et}).
\]
The inverse-image functor $u^{-1}:\Ab(V_{\Et})\to\Ab(U_{\Et})$ is exact,
so it acts termwise on complexes and satisfies
\[
 u^{-1}\dashv Ru_*.
\]
Its unit and counit are denoted $\eta_u$ and $\epsilon_u$.  With a fixed
coefficient ring whose pullback is canonically identified with itself, the
same functors act on module sheaves.
\end{definition}
\deps{Definition 1.1 and the existence of $K$-injective resolutions in a
Grothendieck abelian category.}

\begin{definition}[Cartesian square with chosen identification]
A cartesian square $\square=(f,g,g',f')$ is the data
\[
\begin{CD}
 X' @>{g'}>> X\\
 @V{f'}VV @VV{f}V\\
 Y' @>{g}>> Y
\end{CD}
\]
together with a chosen identification $X'\simeq X\times_YY'$ and the
corresponding equality $fg'=gf'$.  This choice gives a specified natural
isomorphism
\[
 f'^{-1}g^{-1}\xrightarrow{\sim}g'^{-1}f^{-1}.
\]
\end{definition}
\deps{Fiber products of schemes and the functoriality of inverse image.}

\begin{definition}[Underived exchange]
For $F\in\Ab(X_{\Et})$, define
\[
 \Ex^0_{\square,F}:g^{-1}f_*F\longrightarrow f'_*g'^{-1}F
\]
to be the adjoint, under $f'^{-1}\dashv f'_*$, of
\[
 f'^{-1}g^{-1}f_*F\xrightarrow{\sim}
 g'^{-1}f^{-1}f_*F\xrightarrow{g'^{-1}\epsilon_f}g'^{-1}F.
\]
For an etale $V\to Y'$, a representative of a section on the left consists
of an etale $W\to Y$, a map $V\to W\times_YY'$, and a section of $F$ on
$X\times_YW$.  Its image is the restriction of that section along
\[
 X'\times_{Y'}V\longrightarrow X\times_YW,
\]
followed by sheafification.  This sectionwise rule is independent of the
representative.
\end{definition}
\deps{Definitions 1.2 and 1.3 and the unit--counit description of an
adjoint mate.}

\begin{lemma}[Naturality of underived exchange]
For every morphism $a:F\to G$ in $\Ab(X_{\Et})$, the square
\[
\begin{CD}
 g^{-1}f_*F @>{\Ex^0_{\square,F}}>> f'_*g'^{-1}F\\
 @V{g^{-1}f_*a}VV @VV{f'_*g'^{-1}a}V\\
 g^{-1}f_*G @>{\Ex^0_{\square,G}}>> f'_*g'^{-1}G
\end{CD}
\]
commutes.
\end{lemma}
\deps{Definition 1.4 and naturality of the adjunction counit.}

\begin{definition}[Derived exchange]
For $K\in D(X_{\Et})$, the derived exchange map is the adjoint mate
\[
 \Ex_{\square,K}:g^{-1}Rf_*K\longrightarrow Rf'_*g'^{-1}K
\]
of the composite
\[
 f'^{-1}g^{-1}Rf_*K\xrightarrow{\sim}
 g'^{-1}f^{-1}Rf_*K\xrightarrow{g'^{-1}\epsilon_f}g'^{-1}K.
\]
Equivalently, choose $K\to I$ $K$-injective, apply
$\Ex^0_{\square,I}$ term by term, and compare $g'^{-1}I$ with a
$K$-injective resolution on $X'$.  Both constructions define the same
morphism in the derived category.
\end{definition}
\deps{Definitions 1.2--1.4 and the derived adjunction
$f^{-1}\dashv Rf_*$.}

\begin{lemma}[Replacement independence]
Let $K\to I$ and $K\to J$ be two $K$-injective resolutions.  Every
comparison $I\to J$ over $K$ gives the same derived exchange morphism, and
two such comparisons are chain-homotopic.  Consequently
$\Ex_{\square,K}$ is independent of the resolution and is natural in
$K$.
\end{lemma}
\deps{Definition 1.6, the $K$-injective comparison lemma, and homotopy
uniqueness of maps between resolutions.}

\begin{lemma}[Identity and pasting equations]
The transformation $\Ex_{\square}$ is the identity for an identity square.
For horizontally or vertically pasted cartesian squares, the exchange map
of the outer rectangle equals the composite of the exchange maps of the
two inner squares after inserting the canonical associativity
isomorphisms for inverse and direct images.
\end{lemma}
\deps{Definition 1.6, Lemma 1.7, and the triangle identities for the two
adjunctions.}

\begin{definition}[Henselian terminology]
A pair $(A,I)$ consisting of a commutative ring and an ideal is henselian
if, for every finite $A$-algebra $B$, every idempotent of $B/IB$ lifts
uniquely to an idempotent of $B$.  A local ring is henselian when its
maximal-ideal pair is henselian, and it is strictly henselian when it is
henselian and its residue field is separably closed.
\end{definition}
\deps{Finite algebras, idempotents, local rings, and separably closed
fields.}

\begin{definition}[Constructible finite-coefficient complexes]
Let $\Lambda$ be a finite commutative ring.  A sheaf of
$\Lambda_T$-modules on $T_{\Et}$ is lisse if it is locally constant with
finite stalks.  It is constructible if $T$ admits a finite partition by
locally closed constructible subschemes on each of which the restriction is
lisse.  A complex $K\in D(T_{\Et},\Lambda)$ is bounded constructible if
every $\mathcal H^i(K)$ is constructible and there are integers $a\leq b$
such that $\mathcal H^i(K)=0$ for $i\notin[a,b]$.  The corresponding full
triangulated subcategory is denoted
\[
 \Dbc(T_{\Et},\Lambda)\subset D(T_{\Et},\Lambda).
\]
\end{definition}
\deps{Definition 1.1, finite stratifications, and locally constant sheaves
with finite stalks.}

\begin{definition}[Finite coefficients on the pro-etale site]
Let $T_{\Proet}$ be the small pro-etale site of weakly etale
$U\to T$, with jointly surjective fpqc covering families.  For a finite
discrete ring $\Lambda$, let $\Lambda_{T,\Proet}$ be the sheafification of
\[
 U\longmapsto\operatorname{Cont}(\pi_0(|U|),\Lambda).
\]
Let $\Mod(T_{\Proet},\Lambda)$ be its abelian category of module sheaves and
define
\[
 D(T_{\Proet},\Lambda)=
 \mathsf{Cpx}(\Mod(T_{\Proet},\Lambda))
 [\{\text{quasi-isomorphisms}\}^{-1}].
\]
The subcategories $\Dplus(T_{\Proet},\Lambda)$ and
$\Dbc(T_{\Proet},\Lambda)$ are defined by lower boundedness and by pullback
of the bounded constructible etale objects, respectively.
\end{definition}
\deps{Definition 1.10, the pro-etale topology, sheafification, and Verdier
localization.}

\begin{definition}[Derived operations on ringed pro-etale topoi]
Let $u:U\to V$ be a morphism and let $R_V$ be a coefficient sheaf whose
inverse image is identified with $R_U$.  The module inverse image is
\[
 u^*F=u^{-1}F\otimes_{u^{-1}R_V}R_U.
\]
Using a functorial $K$-flat resolution and a functorial $K$-injective
resolution defines ordinary triangulated functors
\[
 Lu^*:D(V_{\Proet},R_V)\longrightarrow D(U_{\Proet},R_U),
 \qquad
 Ru_*:D(U_{\Proet},R_U)\longrightarrow D(V_{\Proet},R_V),
\]
with $Lu^*\dashv Ru_*$.  For finite constant coefficients the displayed
ring map is an isomorphism and inverse image is exact.
\end{definition}
\deps{Definition 1.11, derived tensor product, and $K$-flat and
$K$-injective resolutions.}

\begin{definition}[Ambient ringed exchange]
For a cartesian square as in Definition 1.3 and a coefficient sheaf $R$,
define the ambient pro-etale exchange transformation
\[
 \Ex^R_{\square}:Lg^*Rf_*\longrightarrow Rf'_*Lg'^*
\]
to be the mate, under $Lf'^*\dashv Rf'_*$, of
\[
 Lf'^*Lg^*Rf_*\xrightarrow{\sim}
 Lg'^*Lf^*Rf_*\xrightarrow{Lg'^*\epsilon_f}Lg'^*.
\]
It is defined before any completeness or constructibility condition.
\end{definition}
\deps{Definitions 1.3 and 1.12 and the unit--counit calculus for an
adjunction.}

\section{S0: comparison predicates and injective acyclicity}

\begin{definition}[Proper-base-change property]
Let $f:X\to Y$ be proper.  Write $\mathrm{PBC}(f)$ if, for every
$g:Y'\to Y$, every associated cartesian square $\square$, and every torsion
sheaf $F\in\Ab(X_{\Et})$, the specified map
\[
 \Ex_{\square,F}:g^{-1}Rf_*F\longrightarrow Rf'_*g'^{-1}F
\]
is an isomorphism in $\Dplus(Y'_{\Et})$.
\end{definition}
\deps{Definitions 1.3 and 1.6.}

\begin{definition}[Injective-acyclicity property]
For a prime $r$, write $\mathrm{IA}_r(f)$ if, after every base change
$g:Y'\to Y$, every injective sheaf $I$ of $\F_r$-modules on $X_{\Et}$
satisfies
\[
 R^qf'_*g'^{-1}I=0\qquad(q>0).
\]
Write $\mathrm{IA}(f)$ if this holds for every prime $r$.
\end{definition}
\deps{Definitions 1.2 and 1.3.}

\begin{lemma}[Finite refinement of a surjective etale morphism]
Let $u:U\to T$ be a surjective etale morphism and assume that $T$ is
quasi-compact and quasi-separated.  There exist a finite set $J$, a finite,
finitely presented, surjective morphism
\[
 r:T'\longrightarrow T,
\]
a finite open covering $T'=\bigcup_{a\in J}W_a$, and $T$-morphisms
$\phi_a:W_a\to U$ such that
\[
 u\circ\phi_a=r|_{W_a}:W_a\longrightarrow T
 \qquad(a\in J).
\]
Equivalently, $r$ factors through $u$ Zariski locally on $T'$.
\end{lemma}
\deps{Finite affine refinement of a quasi-compact etale cover, Zariski's
Main Theorem, induction on the bound for its geometric fibers,
approximation by finite finitely presented morphisms, and descent of
quasi-compact opens and morphisms of finite presentation.}
\sources{Stacks Project, More on Morphisms, Lemma 37.45.7, Tag 09Z0;
internal inputs Tags 01KN, 03JA, 02LR, 07RY, 09YZ, 01Z4, and 01ZC.}

\begin{lemma}[Degree-zero comparison over a strictly henselian base]
Let $A$ be strictly henselian local, let $S=\Spec A$, let $s$ be its closed
geometric point, and let $p:Z\to S$ be proper.  For every etale sheaf $F$
on $Z$, restriction gives
\[
 \Gamma(Z,F)\xrightarrow{\sim}\Gamma(Z_s,F|_{Z_s}).
\]
This isomorphism is natural in $F$ and in finite pullback over $S$.
\end{lemma}
\deps{Lemma 2.3 and the henselian degree-zero comparison isolated in
Section 6 below.}
\sources{Stacks Project, Tags 0A0C and 0A3S.}

\begin{lemma}[Underived proper exchange]
For a proper cartesian square and any sheaf $F$ on $X_{\Et}$, the map
\[
 \Ex^0_{\square,F}:g^{-1}f_*F\xrightarrow{\sim}f'_*g'^{-1}F
\]
is an isomorphism.
\end{lemma}
\deps{Definition 1.4, Lemma 2.4, continuity over pointed affine etale
neighborhoods, properness under base change, and conservativity of
geometric stalks.}
\sources{Stacks Project, Tag 0A3U.}

\begin{lemma}[Filtered torsion union]
For every torsion sheaf $F$ on $X_{\Et}$,
\[
 F=\colim_{(\mathbf N_{>0},\mid)}F[n],\qquad
 F[n]=\ker(n:F\to F),
\]
as a filtered colimit in $\Ab(X_{\Et})$.  Here positive integers are
ordered by divisibility, and for $n\mid m$ the transition map is the
inclusion $F[n]\hookrightarrow F[m]$; this poset is filtered because
$\operatorname{lcm}(n,m)$ is a common upper bound.
\end{lemma}
\deps{The definition of a torsion sheaf and exact filtered colimits in a
Grothendieck abelian category.}

\begin{lemma}[Filtered colimits and higher direct image]
If $u:U\to V$ is quasi-compact and quasi-separated and $(F_i)$ is a
filtered system of torsion etale sheaves, then for every $q\geq0$ the
canonical map
\[
 \colim_i R^qu_*F_i\xrightarrow{\sim}R^qu_*(\colim_iF_i)
\]
is an isomorphism.  The same assertion holds after arbitrary base change.
\end{lemma}
\deps{Quasi-compact quasi-separated etale cohomology and exactness of
filtered colimits.}
\sources{Stacks Project, Tag 0A4B.}

\begin{lemma}[Prime-factor devissage]
Let $n>1$, choose a prime $r\mid n$, and let $F$ be killed by $n$.  There
is a natural exact sequence
\[
 0\longrightarrow F[r]\longrightarrow F
 \xrightarrow{r} rF\longrightarrow0,
\]
where $F[r]$ is killed by $r$ and $rF$ is killed by $n/r$.  Thus a property
of torsion sheaves stable under short exact sequences and known for
prime-torsion sheaves holds for every sheaf killed by a positive integer.
\end{lemma}
\deps{Kernels and images in the abelian category of etale sheaves.}

\begin{lemma}[Five-lemma transfer for exchange]
For a short exact sequence $0\to F_1\to F\to F_2\to0$ of torsion sheaves,
the exchange transformations form a morphism between the two long exact
sequences of higher direct images.  If exchange is an isomorphism for
$F_1$ and $F_2$ in every degree, then it is an isomorphism for $F$ in every
degree.
\end{lemma}
\deps{Lemma 1.5, the natural long exact sequence of a right-derived
functor, and the five lemma.}

\begin{lemma}[Leray acyclicity criterion]
Let $F\to I^\bullet$ be an injective resolution and let $v$ be a left exact
functor.  If every $I^m$ is $v$-acyclic, then $Rv(F)$ is represented by
$v(I^\bullet)$.  Applied to $v=f'_*$ and $g'^{-1}I^m$, this identifies the
cohomology of the target of exchange with the cohomology of the termwise
underived target.
\end{lemma}
\deps{Injective resolutions and the acyclic-resolution lemma for derived
functors.}

\begin{proposition}[Equivalence of proper base change and acyclicity]
For a proper morphism $f$, the following conditions are equivalent.
\begin{enumerate}
\item $\mathrm{PBC}(f)$ holds.
\item $\mathrm{IA}_r(f)$ holds for every prime $r$.
\end{enumerate}
\end{proposition}
\deps{Definitions 2.1 and 2.2 and Lemmas 2.5--2.10.}
\sources{Stacks Project, Tags 0A4A and 0A4B.}

\section{S1: finite morphisms}

\begin{lemma}[Geometric stalk formula]
Let $f:X\to Y$ be finite, let $F$ be an etale sheaf of abelian groups on
$X$, and let $\bar y\to Y$ be a geometric point.  If
$\{\bar x\}$ is the finite set of geometric points of $X_{\bar y}$, then
\[
 (f_*F)_{\bar y}\xrightarrow{\sim}
 \prod_{\bar x\in|X_{\bar y}|}F_{\bar x}.
\]
For a finite set, this product is also the direct sum.
\end{lemma}
\deps{The definition of a geometric stalk and the decomposition of a
finite scheme over a strictly henselian local scheme into local factors.}
\sources{Stacks Project, Proposition 59.55.2, Tag 03QP.}

\begin{lemma}[Exactness of finite direct image]
For finite $f$, the functor
$f_*:\Ab(X_{\Et})\to\Ab(Y_{\Et})$ is exact.  Consequently
$R^qf_*=0$ for $q>0$ and $Rf_*=f_*$ on $D^+(X_{\Et})$.
\end{lemma}
\deps{Lemma 3.1, exactness of finite products, and conservativity of
geometric stalks.}

\begin{lemma}[Finite exchange]
In a cartesian square with $f$ finite, the underived exchange map is the
bijection obtained by base changing the finite geometric fiber:
\[
 g^{-1}f_*F\xrightarrow{\sim}f'_*g'^{-1}F.
\]
The isomorphism is compatible with the stalk formula on both sides.
\end{lemma}
\deps{Definition 1.4 and Lemma 3.1.}

\begin{proposition}[Finite morphisms satisfy proper base change]
Every finite morphism $f$ satisfies $\mathrm{PBC}(f)$.
\end{proposition}
\deps{Lemmas 3.2 and 3.3, or equivalently Proposition 2.11.}
\sources{Stacks Project, Tag 0A4E.}

\section{S2: composition and proper-surjective cancellation}

\begin{lemma}[Direct image preserves injectives]
For every morphism $u:U\to V$, the functor $u_*$ sends injective etale
sheaves to injective etale sheaves.  The relevant adjunction is
$u^{-1}\dashv u_*$, and $u^{-1}$ is exact.
\end{lemma}
\deps{Definition 1.2 and the elementary injective-object criterion for a
right adjoint to an exact functor.}

\begin{lemma}[Grothendieck spectral sequence for a composite]
Let $X\xrightarrow{f}Y\xrightarrow{h}Z$ be morphisms and let $F$ be a
sheaf such that $f_*F$ is $h_*$-acyclic, as happens when $F$ is injective.
There is a functorial first-quadrant spectral sequence
\[
 E_2^{p,q}=R^ph_*R^qf_*F
 \Longrightarrow R^{p+q}(hf)_*F.
\]
It is compatible with inverse image in a cartesian base change.
\end{lemma}
\deps{Lemma 4.1 and the Grothendieck spectral-sequence theorem for two
left exact functors.}

\begin{proposition}[Composition stability]
Let $X\xrightarrow{f}Y\xrightarrow{h}Z$ be proper.  If
$\mathrm{IA}(f)$ and $\mathrm{IA}(h)$ hold, then
$\mathrm{IA}(hf)$ and hence $\mathrm{PBC}(hf)$ hold.
\end{proposition}
\deps{Proposition 2.11 and Lemmas 2.5, 4.1, and 4.2.}
\sources{Stacks Project, Tag 0A4D.}

\begin{lemma}[Retract supplied by a proper surjection]
Let $f:X\to Y$ be proper and surjective, let $I$ be injective on
$Y_{\Et}$, and choose a monomorphism $f^{-1}I\hookrightarrow J$ with $J$
injective on $X_{\Et}$.  The adjoint map
\[
 I\longrightarrow f_*J
\]
is a monomorphism.  Injectivity of $I$ supplies a retraction
$f_*J\to I$, so $I$ is a direct summand of $f_*J$.
\end{lemma}
\deps{Surjectivity on geometric points, the stalk definition of direct
image, and injectivity of $I$.}

\begin{lemma}[Vanishing descends to a retract]
If $I$ is a direct summand of $M$ and $R^qh_*M=0$ for $q>0$, then
$R^qh_*I=0$ for $q>0$.  This remains true after arbitrary base change.
\end{lemma}
\deps{Additivity of derived functors and preservation of retracts by an
additive functor.}

\begin{proposition}[Cancellation along a proper surjection]
Let $X\xrightarrow{f}Y\xrightarrow{h}Z$ be proper, with $f$ surjective.
If $\mathrm{PBC}(f)$ and $\mathrm{PBC}(hf)$ hold, then
$\mathrm{PBC}(h)$ holds.
\end{proposition}
\deps{Proposition 2.11, Lemmas 4.2, 4.4, and 4.5, and the underived
proper exchange of Lemma 2.5.}
\sources{Stacks Project, Tag 0A4C.}

\section{S3: reduction to the projective line}

\begin{lemma}[Target locality]
For proper $f:X\to Y$, the assertion $\mathrm{PBC}(f)$ is local on $Y$ for
the etale topology.  More precisely, after an etale covering
$\coprod_iV_i\to Y$, the exchange map is an isomorphism if and only if its
pullback over every $V_i$ is an isomorphism.
\end{lemma}
\deps{The description of $R^qf_*F$ as the sheafification of
$V\mapsto H^q(X\times_YV,F|)$ and conservativity of geometric stalks.}
\sources{Stacks Project, Tag 03Q8.}

\begin{proposition}[Chow and projective factorization]
Let $f:X\to Y$ be proper and let $Y$ be affine.  There exist a proper
surjection $\pi:\widetilde X\to X$, an integer $n\geq0$, and a closed
immersion
\[
 i_Y:\widetilde X\hookrightarrow\mathbf P^n_Y
\]
such that $f\pi$ is the composite of $i_Y$ with
$\mathbf P^n_Y\to Y$.  The assertion is stable under base change on $Y$.
\end{proposition}
\deps{Lemma 5.1, approximation by finitely generated rings, noetherian
Chow's lemma, and the fact that a proper immersion is closed.}
\sources{Stacks Project, Tag 0202.}

\begin{lemma}[Projective embedding over the source]
With the notation of Proposition 5.2, the map
\[
 i_X=(\pi,i_Y):\widetilde X\longrightarrow
 X\times_Y\mathbf P^n_Y\simeq\mathbf P^n_X
\]
is a closed immersion over $X$.
\end{lemma}
\deps{Proposition 5.2, the graph criterion for properness, and stability of
immersions under base change.}

\begin{definition}[Symmetric-product coordinate morphism]
For $n\geq1$ and a scheme $S$, define polynomials $F_0,\dots,F_n$ by
\[
 \prod_{i=1}^n(x_it+y_i)=\sum_{j=0}^nF_jt^j.
\]
Each $F_j$ is multihomogeneous of degree one in every pair $(x_i,y_i)$.
The associated rational formula is
\[
 ((x_1:y_1),\dots,(x_n:y_n))\longmapsto[F_0:\cdots:F_n].
\]
\end{definition}
\deps{The universal expansion of a product of linear binary forms.}

\begin{lemma}[No common zero and permutation invariance]
The polynomials $F_0,\dots,F_n$ have no common zero on
$(\mathbf P^1_S)^n$ and are invariant under the action of the symmetric
group $S_n$.  Hence Definition 5.4 gives an $S_n$-invariant morphism
\[
 \sigma_n:(\mathbf P^1_S)^n\longrightarrow\mathbf P^n_S.
\]
\end{lemma}
\deps{Definition 5.4 and the fact that a product of nonzero linear forms
is a nonzero homogeneous polynomial.}

\begin{lemma}[Finiteness of the symmetric-product morphism]
The morphism $\sigma_n$ is finite.  More precisely, on the chart where all
$y_i$ are nonzero, the induced affine ring extension is integral: each
$z_i=x_i/y_i$ satisfies
\[
 T^n-e_1T^{n-1}+\cdots+(-1)^ne_n=0.
\]
The corresponding reversed equations give integral affine ring extensions
on the charts at infinity, and these charts cover $(\mathbf P^1_S)^n$.
\end{lemma}
\deps{Lemma 5.5, the affine criterion for finite morphisms, and integrality
of roots over the elementary-symmetric subring.}

\begin{lemma}[Surjectivity of the symmetric-product morphism]
The finite morphism $\sigma_n$ is surjective.  Equivalently, over every
algebraically closed geometric point of $S$, every binary form of degree
$n$ factors as a product of $n$ linear forms, including factors at
infinity.
\end{lemma}
\deps{Lemma 5.6, factorization of one-variable polynomials over an
algebraically closed field, and the geometric criterion for surjectivity.}

\begin{lemma}[Reduction of projective space to products of lines]
If every projection $\mathbf P^1_T\to T$ satisfies proper base change,
then $(\mathbf P^1_T)^n\to T$ does so by repeated base change and
composition, and $\mathbf P^n_T\to T$ does so by cancellation along
$\sigma_n$.
\end{lemma}
\deps{Propositions 3.4, 4.3, and 4.6 and Lemmas 5.6--5.7.}

\begin{proposition}[Reduction to the projective line]
Assume that $\mathbf P^1_S\to S$ satisfies $\mathrm{PBC}$ for every
scheme $S$.  Then every proper morphism of schemes satisfies
$\mathrm{PBC}$.
\end{proposition}
\deps{Lemma 5.1, Proposition 5.2, Lemmas 5.3 and 5.8, finite-morphism base
change, composition stability, and proper-surjective cancellation.}
\sources{Stacks Project, Tag 0A4F.}

\section{S4: the projective line over a strictly henselian base}

\subsection{Henselian pairs and degree zero}

\begin{definition}[Henselian pair]
A pair $(A,I)$ consisting of a ring and an ideal is henselian if, for every
finite $A$-algebra $B$, every idempotent of $B/IB$ lifts uniquely to an
idempotent of $B$.  For $T=\Spec A$ put $T_0=\Spec(A/I)$.  A strictly
henselian local ring is a henselian local ring with separably closed
residue field.
\end{definition}
\deps{Finite algebras and idempotents.}

\begin{definition}[Strictly henselian projective-line setup]
Let $(A,\mathfrak m)$ be strictly henselian local, put
$S=\Spec A$, $s=\Spec(A/\mathfrak m)$,
$X=\mathbf P^1_S$, and $X_0=\mathbf P^1_s$.  For an etale sheaf $F$ on
$X$, write $F_0$ for its inverse image on $X_0$ and let
\[
 \operatorname{res}_F:\Gamma(X,F)\longrightarrow\Gamma(X_0,F_0)
\]
be restriction.
\end{definition}
\deps{Definition 6.1 and inverse image along the closed immersion
$X_0\hookrightarrow X$.}

\begin{lemma}[Restriction on a spectral space]
Let $T$ be spectral and $T_0\subset T$ closed.  Suppose every irreducible
closed subset $C\subset T$ meets $T_0$ in a nonempty connected subset.
Then restriction is a bijection
\[
 \Gamma(T,F)\xrightarrow{\sim}\Gamma(T_0,F|_{T_0})
\]
for every sheaf of abelian groups on the Zariski space $T$.
\end{lemma}
\deps{The specialization order on a spectral space and the sheaf gluing
axiom.}
\sources{Stacks Project, Tag 09ZG.}

\begin{proposition}[Stein connectedness over a henselian pair]
Let $(A,I)$ be henselian and let $p:Y\to\Spec A$ be proper.  In the Stein
factorization
\[
 Y\longrightarrow\Spec B\longrightarrow\Spec A,
\]
$A\to B$ is finite and $(B,IB)$ is henselian.  Every irreducible closed
subset of $Y$ meets $Y_0$ in a nonempty connected subset.  Therefore the
Zariski restriction criterion of Lemma 6.3 applies to $(Y,Y_0)$.
\end{proposition}
\deps{Definition 6.1, Stein factorization, connectedness of geometrically
connected fibers, and the closed-surjection connected-fiber criterion.}
\sources{Stacks Project, Tag 0A0B.}

\begin{definition}[Finite pullback of a section]
For a finite morphism $h:V\to Y$ and $a\in\Gamma(Y,F)$, define
$h^\sharp a\in\Gamma(V,h^{-1}F)$ by the inverse-image unit.  Its closed
fiber is the pullback of $a|_{Y_0}$ along $h_0:V_0\to Y_0$.
\end{definition}
\deps{The inverse-image functor and its unit.}

\begin{lemma}[Finite-pullback identities]
Finite pullback of sections satisfies the following four equations:
identity pullback fixes a section; pullback along a composite equals the
composite pullback; pullback commutes with restriction to the closed fiber;
and a section descended from a finite surjective refinement pulls back to
the chosen local section on that refinement.
\end{lemma}
\deps{Definition 6.5, functoriality of inverse image, and uniqueness in
finite descent.}
\sources{Stacks Project, Tag 09ZF, finite-pullback compatibility.}

\begin{proposition}[Etale degree-zero henselian comparison]
Let $(A,I)$ be henselian and let $Y\to\Spec A$ be proper.  For every etale
abelian sheaf $F$ on $Y$,
\[
 \Gamma(Y,F)\xrightarrow{\sim}\Gamma(Y_0,F|_{Y_0}).
\]
The isomorphism is natural in $F$, in finite pullback, and in morphisms of
henselian pairs.
\end{proposition}
\deps{Lemma 2.3, Lemma 6.3, Proposition 6.4, and Lemma 6.6.}
\sources{Stacks Project, Tags 09ZF, 0A0C, and 0A3S.}

\subsection{Gabber's affine comparison}

\begin{lemma}[Finite-stage killing]
Let $(T,T_0)$ be an affine henselian pair, let $F$ be torsion, and let
$c\in H^q(T_0,F|_{T_0})$ with $q>0$.  Suppose $c$ vanishes after pullback
to an integral strictly local cover that is the inverse limit of finite,
finitely presented covers $T_\lambda\to T$.  Then $c$ already vanishes
after pullback to one $T_{\lambda_0}$.
\end{lemma}
\deps{Finite presentation of a cocycle and its null-homotopy and
continuity over a filtered inverse system.}
\sources{Stacks Project, Tag 09ZD.}

\begin{lemma}[Constructible approximation of one class]
Under the hypotheses of Lemma 6.8, a cocycle representing $c$ factors,
after a quasi-compact etale neighborhood, through a monomorphism
$C\hookrightarrow F$ with $C$ a constructible torsion sheaf.
\end{lemma}
\deps{Finite split hypercover skeleta, finitely many section relations,
and geometric-stalk detection of monomorphisms.}

\begin{lemma}[Pushout monomorphism]
If $A\hookrightarrow B$ and $A\hookrightarrow C$ are monomorphisms of
etale sheaves, then the canonical map
\[
 C\longrightarrow B\amalg_AC
\]
is a monomorphism.
\end{lemma}
\deps{The sheaf cokernel description of a pushout in an abelian category.}

\begin{lemma}[Embedding a constructible sheaf in finite pushforwards]
If $U$ is quasi-compact and quasi-separated and $C$ is constructible
torsion on $U_{\Et}$, there is a monomorphism
\[
 C\hookrightarrow\bigoplus_{i=1}^m(h_i)_*\underline{M_i},
\]
where every $h_i:V_i\to U$ is finite and finitely presented and every
$M_i$ is a finite abelian group.
\end{lemma}
\deps{Constructible stratification, noetherian induction after
approximation, and the finite stalk formula of Lemma 3.1.}
\sources{Stacks Project, Tag 09Z7.}

\begin{lemma}[Finite-cover killing with a monic unit]
For a constructible $C$ and $c\in H^q(T_0,C_0)$, $q>0$, there is a finite
finitely presented cover $h:V\to T$ such that $h_0^{-1}c=0$ and the unit
\[
 C\longrightarrow h_*h^{-1}C
\]
is a monomorphism.
\end{lemma}
\deps{Lemmas 3.1, 6.8, and 6.11.}

\begin{proposition}[Gabber effacement]
For an affine henselian pair $(T,T_0)$, a torsion sheaf $F$ on $T_{\Et}$,
and $c\in H^q(T_0,F_0)$ with $q>0$, there is a monomorphism
$F\hookrightarrow F'$ such that the image of $c$ in
$H^q(T_0,F'_0)$ is zero.
\end{proposition}
\deps{Lemmas 6.9--6.12 and a pushout of the finite-cover monomorphism along
the constructible approximation.}
\sources{Stacks Project, Tag 09ZE.}

\begin{proposition}[Gabber affine henselian-pair comparison]
For every affine henselian pair $(T,T_0)$ and every torsion etale sheaf
$F$ on $T$,
\[
 H^q(T,F)\xrightarrow{\sim}H^q(T_0,F_0)
 \qquad(q\geq0).
\]
\end{proposition}
\deps{Propositions 6.7 and 6.13 and induction on $q$ in the two long exact
sequences.}
\sources{Stacks Project, Tag 09ZI.}

\begin{lemma}[Acyclicity bound for restricted injectives]
Let $Y$ be separated with affine diagonal and covered by $n+1$ affine
opens.  If $I$ is injective on $Y_{\Et}$ and $Y_0$ is the closed member of
a henselian pair, then
\[
 H^q(Y_0,I|_{Y_0})=0\qquad(q>n).
\]
\end{lemma}
\deps{Proposition 6.14 and induction on the number of affine opens using
the Mayer--Vietoris long exact sequence.}
\sources{Stacks Project, Tag 0A51.}

\subsection{The remaining degree-one class}

\begin{definition}[Finite torsors]
For a finite abelian group $M$, an $M$-torsor on $Y_{\Et}$ is a finite
locally constant sheaf $P$ with an $M$-action such that
$M\times P\to P\times_YP$ is an isomorphism.  The pointed set of
isomorphism classes is canonically $H^1(Y,\underline M)$.
\end{definition}
\deps{Nonabelian degree-one cohomology and finite locally constant
sheaves.}

\begin{lemma}[Finite locally constant sheaves and finite etale schemes]
Finite locally constant sheaves of sets on $Y_{\Et}$ are equivalent to
finite etale $Y$-schemes.  Under this equivalence, an action of a finite
group is an action by finite-etale automorphisms, and torsors correspond to
simply transitive actions.
\end{lemma}
\deps{Finite etale descent and the sheaf of sections of a finite etale
scheme.}

\begin{proposition}[Finite-etale lifting over a proper henselian pair]
Let $(A,I)$ be henselian, let $Y\to\Spec A$ be proper, and let
$Y_0=Y\times_AA/I$.  Restriction is an equivalence
\[
 \mathsf{FEt}(Y)\xrightarrow{\sim}\mathsf{FEt}(Y_0).
\]
It lifts morphisms and finite group actions uniquely.
\end{proposition}
\deps{Definition 6.1, finite-etale formal lifting, proper algebraization,
and approximation from the noetherian complete case.}
\sources{Stacks Project, Tags 0A48 and 0A5G.}

\begin{corollary}[Torsor comparison]
Under the hypotheses of Proposition 6.18, restriction is a bijection
\[
 H^1(Y,\underline M)\xrightarrow{\sim}H^1(Y_0,\underline M)
\]
for every finite abelian group $M$.
\end{corollary}
\deps{Definition 6.16, Lemma 6.17, and Proposition 6.18.}

\begin{lemma}[Finite-pushforward representation of an $H^1$ class]
Let $I$ be an injective sheaf of $\F_r$-modules on
$X=\mathbf P^1_A$.  Every class $c\in H^1(X_0,I_0)$ is the image of a
finite tuple of classes belonging to sheaves
\[
 (h_i)_*\underline{M_i},
\]
where $h_i:V_i\to X$ is finite and finitely presented and $M_i$ is a
finite $\F_r$-module.
\end{lemma}
\deps{Lemmas 6.9 and 6.11 and injectivity of $I$ to extend the resulting
map of constructible sheaves.}

\begin{proposition}[Acyclicity on the closed projective line]
For an injective sheaf $I$ of $\F_r$-modules on
$X=\mathbf P^1_A$,
\[
 H^q(X_0,I_0)=0\qquad(q>0).
\]
\end{proposition}
\deps{Lemma 6.15, Corollary 6.19, Lemma 6.20, and exactness of finite
direct image.}
\sources{Stacks Project, Tag 0A5H.}

\subsection{Strictly henselian projective-line comparison}

\begin{proposition}[Cohomology comparison on the projective line]
In the setup of Definition 6.2, for every torsion etale sheaf $F$ on $X$,
restriction gives natural isomorphisms
\[
 H^q(X,F)\xrightarrow{\sim}H^q(X_0,F_0)
 \qquad(q\geq0).
\]
\end{proposition}
\deps{Proposition 6.7, Proposition 6.21, an injective resolution of $F$,
and the acyclic-resolution lemma.}
\sources{Stacks Project, Tags 0A0C and 0A3S.}

\begin{proposition}[Strict-local projective-line exchange]
Let $A\to B$ be a local morphism of strictly henselian local rings
inducing a specified extension of their separably closed residue fields.
For the base-change square of $\mathbf P^1_A\to\Spec A$, the stalk of the
exchange map at the closed geometric point of $\Spec B$ is the restriction
map between the cohomologies of the two closed projective lines.  If the
residue-field map is an isomorphism, this stalk map is an isomorphism.
\end{proposition}
\deps{Definition 1.6 and Proposition 6.22, applied termwise to an injective
resolution.}

\section{S5: the projective line over an arbitrary base}

\subsection{Strict-henselization stalks}

\begin{definition}[Pointed etale neighborhoods]
For a geometric point $\bar s:\Spec\Omega\to S$, let
$\mathcal N_{\bar s}$ be the category whose objects are pairs
$(V,\bar v)$ with $V\to S$ etale and $\bar v\to V$ above $\bar s$.
A morphism preserves the chosen geometric points.  Stalk colimits are
indexed by $\mathcal N_{\bar s}^{\mathrm{op}}$.
\end{definition}
\deps{Geometric points and pointed etale morphisms.}

\begin{lemma}[Affine neighborhoods are cofinal]
The full subcategory of affine pointed etale neighborhoods is cofinal for
the opposite-indexed stalk colimit.
\end{lemma}
\deps{Every point of a scheme has an affine open neighborhood and finite
fiber products of pointed neighborhoods.}

\begin{definition}[Strict henselization as a limit]
For affine $(V,\bar v)$ with $V=\Spec A_V$, put
\[
 \mathcal O^{\mathrm{sh}}_{S,\bar s}
   =\colim_{(V,\bar v)}A_V,
 \qquad
 S_{(\bar s)}=\Spec\mathcal O^{\mathrm{sh}}_{S,\bar s}.
\]
Thus $S_{(\bar s)}$ is the inverse limit of the affine pointed etale
neighborhoods.
\end{definition}
\deps{Definitions 7.1, Lemma 7.2, and the anti-equivalence between affine
schemes and rings.}

\begin{lemma}[Fiber products commute with the neighborhood limit]
If $f:X\to S$ is quasi-compact and quasi-separated, then
\[
 X\times_SS_{(\bar s)}\xrightarrow{\sim}
 \varprojlim_{(V,\bar v)}X\times_SV.
\]
\end{lemma}
\deps{Definition 7.3, commutation of tensor products with filtered
colimits, and gluing over a finite affine cover.}

\begin{lemma}[Finite hypercover skeleta descend]
Let $(X_i)$ be a cofiltered inverse system of quasi-compact
quasi-separated schemes with affine transition maps and limit $X$.  Every
finite skeleton of a split etale hypercover of $X$, every morphism between
two such skeleta, and every equality between such morphisms descend to one
finite stage $X_i$.
\end{lemma}
\deps{Finite presentation of etale morphisms, quasi-compactness of matching
objects, and descent of finitely many simplicial identities.}

\begin{proposition}[Etale cohomology continuity]
Under the hypotheses of Lemma 7.5, let a torsion sheaf $F_i$ on $X_i$
pull back to $F$ on $X$.  Then
\[
 \colim_{j\geq i}H^q(X_j,F_j)
 \xrightarrow{\sim}H^q(X,F)
 \qquad(q\geq0).
\]
\end{proposition}
\deps{Lemma 7.5 and the computation of etale cohomology by finite split
hypercover skeleta.}
\sources{Stacks Project, Tag 03Q6.}

\begin{proposition}[Strict-henselization stalk formula]
Let $f:X\to S$ be quasi-compact and quasi-separated, let $F$ be a torsion
etale sheaf, and let $\bar s\to S$ be geometric.  Then
\[
 (R^qf_*F)_{\bar s}\xrightarrow{\sim}
 H^q(X\times_SS_{(\bar s)},F|)
 \qquad(q\geq0).
\]
For $K\in\Dplus(X_{\Et})$ with torsion cohomology, this assembles to
\[
 (Rf_*K)_{\bar s}\xrightarrow{\sim}
 \RGamma(X\times_SS_{(\bar s)},K|)
 \quad\text{in }\Dplus(\Ab).
\]
\end{proposition}
\deps{Definitions 7.1 and 7.3, Lemmas 7.2 and 7.4, Proposition 7.6, and a
bounded-below injective resolution.}
\sources{Stacks Project, Tag 03Q9.}

\subsection{Proper curves over separably closed extensions}

\begin{proposition}[Universal-homeomorphism invariance]
If $u:X'\to X$ is a universal homeomorphism, base change gives an
equivalence of etale sites
\[
 X_{\Et}\xrightarrow{\sim}X'_{\Et}.
\]
The equivalence preserves and reflects covering families and is compatible
with identities, composition, and pullback of sheaves and cohomology.
\end{proposition}
\deps{Effective descent for etale morphisms across an integral radicial
surjection and nilpotent invariance of etale algebras.}
\sources{Stacks Project, Tag 03SI, including its covering compatibility.}

\begin{definition}[Eight curve assertions]
Let $k$ be algebraically closed, let $C$ be separated of finite type over
$k$ with $\dim C\leq1$, and let $F\in\Ab(C_{\Et})$ be a torsion sheaf.
A torsion sheaf on $C_{\Et}$ is called constructible if $C$ has a finite
partition by locally closed constructible subschemes on each of which its
restriction is locally constant with finite stalks.
When $\operatorname{char}k=p>0$, call $F$ $p$-primary if every element of
every geometric stalk $F_{\bar c}$ is killed by a power of $p$, and call
$F$ prime to $p$ if every such element has finite order coprime to $p$.
In characteristic zero, ``prime to the characteristic'' imposes no
condition and the third assertion below has no instances.

For an algebraically closed extension $k'/k$, write
\[
 \pi_{k'}:C_{k'}=C\times_k k'\longrightarrow C,
 \qquad F_{k'}=\pi_{k'}^{-1}F,
\]
and let
\[
 \rho^q_{k'/k}(C,F):H^q(C_{\Et},F)\longrightarrow
 H^q((C_{k'})_{\Et},F_{k'})
\]
be the pullback map.  For $1\leq j\leq8$, the assertion
$\mathcal C_j(C,F)$ is the following fully quantified statement.
\begin{enumerate}
\item[\textnormal{(C1)}] For every integer $q>2$,
$H^q(C_{\Et},F)=0$.
\item[\textnormal{(C2)}] If $C$ is affine, then for every integer $q>1$,
$H^q(C_{\Et},F)=0$.
\item[\textnormal{(C3)}] If $\operatorname{char}k=p>0$ and $F$ is
$p$-primary, then for every integer $q>1$, $H^q(C_{\Et},F)=0$.
\item[\textnormal{(C4)}] If $F$ is constructible and prime to
$\operatorname{char}k$, then $H^q(C_{\Et},F)$ is a finite abelian group
for every integer $q\geq0$.
\item[\textnormal{(C5)}] If $C$ is proper over $k$ and $F$ is
constructible, then $H^q(C_{\Et},F)$ is a finite abelian group for every
integer $q\geq0$.
\item[\textnormal{(C6)}] If $F$ is prime to $\operatorname{char}k$, then
for every algebraically closed extension $k'/k$ and every integer
$q\geq0$, the map $\rho^q_{k'/k}(C,F)$ is an isomorphism.
\item[\textnormal{(C7)}] If $C$ is proper over $k$, then for every
algebraically closed extension $k'/k$ and every integer $q\geq0$, the map
$\rho^q_{k'/k}(C,F)$ is an isomorphism, without a restriction on the
torsion of $F$.
\item[\textnormal{(C8)}] For every open immersion $a:U\hookrightarrow C$,
the restriction map
\[
 a^*:H^2(C_{\Et},F)\longrightarrow H^2(U_{\Et},a^{-1}F)
\]
is surjective.
\end{enumerate}
Write $\mathcal C(C,F)$ for the conjunction
$\bigwedge_{j=1}^8\mathcal C_j(C,F)$.
\end{definition}
\deps{Etale cohomology, geometric stalks, constructible torsion sheaves,
and exact inverse image on the small etale site.}
\sources{Stacks Project, Tag 03SC, numbered assertions (1)--(8).}

\begin{lemma}[Filtered cohomology on a curve]
For a quasi-compact quasi-separated curve $C$ and a filtered system of
torsion etale sheaves $(F_i)$,
\[
 \colim_iH^q(C,F_i)\xrightarrow{\sim}
 H^q(C,\colim_iF_i)
\]
for every $q$.  This isomorphism commutes with extension of the ground
field.
\end{lemma}
\deps{Finite split hypercover skeleta and exactness of filtered colimits.}

\begin{lemma}[Short-exact devissage for the eight assertions]
Let
\[
 0\longrightarrow F_1\longrightarrow F\longrightarrow F_2
 \longrightarrow0
\]
be exact in $\Ab(C_{\Et})$.  If $\mathcal C(C,F_1)$ and
$\mathcal C(C,F_2)$ hold, then $\mathcal C(C,F)$ holds.  In particular,
the class of torsion sheaves satisfying all eight implications of
Definition 7.9 is closed under extensions.
\end{lemma}
\deps{Definition 7.9, exact inverse image, natural long exact cohomology
sequences, the five lemma, the Serre property of constructible torsion
sheaves, and assertion $\mathcal C_1$ for the $\mathcal C_8$ transfer.}
\sources{Stacks Project, Tag 0A5C.}

\begin{lemma}[Reduction to constructible prime torsion]
For a proper finite-type curve $C$ and a torsion sheaf $F$, assertion
$\mathcal C_7(C,F)$ follows from its instances for reduced curves with
constructible $\F_r$-coefficients, for every prime $r$.  The intervening
coefficient passage is governed by constructible subsheaves, finite
primary decomposition, and the exact sequence
\[
 0\to F[r]\to F\to F/F[r]\to0.
\]
\end{lemma}
\deps{Proposition 7.8 and Lemmas 7.10 and 7.11.}

\begin{lemma}[Dense-open localization]
Let $C$ be a reduced curve and $F$ a constructible $\F_r$-sheaf.  There is
a dense open $j:U\hookrightarrow C$, a finite disjoint union of irreducible
curves, on which $F$ is finite locally constant.  For
$i:Z=C\setminus U\hookrightarrow C$ there is a natural exact sequence
\[
 0\to j_!j^{-1}F\to F\to i_*i^{-1}F\to0,
\]
and $Z$ is zero-dimensional.  Write $U=\coprod_aU_a$ and let
$j_a:U_a\hookrightarrow C$ be the composite open immersion.  Then
\[
 \mathcal C(C,F)\quad\Longleftrightarrow\quad
 \mathcal C(C,j_{a,!}(F|_{U_a}))\text{ for every }a.
\]
\end{lemma}
\deps{Constructible stratification, geometric-stalk exactness, and Lemma
7.11.}

\begin{lemma}[Finite $r$-group filtration]
Let $H$ be a finite $r$-group and let $V$ be a nonzero finite-dimensional
$\F_r[H]$-module.  Then $V^H\neq0$, and $V$ has an $H$-stable finite
filtration whose successive quotients are the trivial one-dimensional
representation.
\end{lemma}
\deps{Nilpotence of the augmentation ideal of $\F_r[H]$ and induction on
$\dim_{\F_r}V$.}
\sources{Stacks Project, Tag 0A3R, finite $r$-group representation step.}

\begin{lemma}[Locally constant reduction and trace splitting]
Let $U$ be irreducible and let $G$ be a finite locally constant
$\F_r$-sheaf.  There is a connected finite etale map $v:V\to U$ of degree
prime to $r$ such that $v^{-1}G$ has a filtration with constant
$\F_r$-quotients.  The adjunction unit and trace satisfy
\[
 G\xrightarrow{\eta}v_*v^{-1}G
 \xrightarrow{\operatorname{tr}}G
   =\deg(v)\operatorname{id}_G,
\]
so $G$ is a direct summand of $v_*v^{-1}G$.
\end{lemma}
\deps{Lemma 7.14, a Sylow subgroup of the finite monodromy group, and the
finite geometric-stalk formula.}
\sources{Stacks Project, Tag 0A3R.}

\begin{definition}[Finite pushforward and extension by zero]
For a cartesian square
\[
\begin{CD}
 U @>{j_U}>> X\\
 @V{v}VV @VV{\bar v}V\\
 V @>{j_V}>> Y,
\end{CD}
\]
with $j_U,j_V$ open immersions and $\bar v$ finite, define
\[
 \beta_{\bar v,j_V,G}:(j_V)_!v_*G
 \longrightarrow\bar v_*(j_U)_!G
\]
as the adjoint of the finite-base-change identification after restriction
to $V$.
\end{definition}
\deps{The adjunction $j_!\dashv j^{-1}$ and finite underived exchange.}

\begin{proposition}[Finite pushforward commutes with extension by zero]
The morphism $\beta_{\bar v,j_V,G}$ of Definition 7.16 is an isomorphism.
It is natural in $G$, compatible with identity and composite finite maps,
and stable under extension of the ground field.
\end{proposition}
\deps{Definition 7.16, the sectionwise support condition for $j_!$, and the
finite stalk formula.}
\sources{Stacks Project, Tag 0GJ1, including functoriality.}

\begin{proposition}[Smooth-curve calculations]
Let $k$ be algebraically closed and let $C$ be a smooth separated curve.
If $r\neq\operatorname{char}k$, the constant sheaf $\F_r$ satisfies
$\mathcal C_1$, $\mathcal C_4$, $\mathcal C_6$, and $\mathcal C_8$; it
also satisfies $\mathcal C_2$ when $C$ is affine and
$\mathcal C_5,\mathcal C_7$ when $C$ is proper.  If
$\operatorname{char}k=p>0$, the constant sheaf $\F_p$ satisfies
$\mathcal C_1$, $\mathcal C_3$, and $\mathcal C_8$; it also satisfies
$\mathcal C_2$ when $C$ is affine and $\mathcal C_5,\mathcal C_7$ when
$C$ is proper.  The maps in $\mathcal C_6$, $\mathcal C_7$, and
$\mathcal C_8$ are the pullback and restriction maps fixed in Definition
7.9.
\end{proposition}
\deps{Kummer theory, the Jacobian and degree sequence, localization at the
finite boundary, the Artin--Schreier sequence, and coherent proper base
change for $H^1(C,\mathcal O_C)$.}
\sources{Stacks Project, Tags 03SG and 0A5B.}

\begin{proposition}[Normalization comparison]
Let $C$ be a reduced curve, let $\nu:\widetilde C\to C$ be its finite
normalization, and let $j:U\hookrightarrow C$ be open.  There is a natural
exact sequence
\[
 0\to j_!\F_r\longrightarrow
 \nu_*\widetilde j_!\F_r\longrightarrow Q\to0,
\]
where $Q$ is constructible and supported on the finite nonnormal locus.
The sequence, its boundary maps, and its restriction to opens commute with
algebraically closed field extension.
\end{proposition}
\deps{Proposition 7.17 and the diagonal finite-fiber map for the
normalization.}
\sources{Stacks Project, Tag 0A3Q.}

\begin{proposition}[Algebraically closed extension for proper curves]
Let $k\subset k'$ be algebraically closed, let $C$ be proper of dimension
at most one over $k$, and let $F$ be a torsion etale sheaf.  Then
\[
 H^q(C,F)\xrightarrow{\sim}H^q(C_{k'},F|_{C_{k'}})
 \qquad(q\geq0).
\]
\end{proposition}
\deps{Lemmas 7.12--7.15 and Propositions 7.17--7.19.}
\sources{Stacks Project, Tags 03SC and 0A5E.}

\begin{proposition}[Separably closed extension for proper curves]
Let $k\subset k'$ be separably closed, let $C$ be proper of dimension at
most one over $k$, and let $F$ be torsion.  The same canonical restriction
map is an isomorphism in every degree.
\end{proposition}
\deps{Proposition 7.20, passage to compatible algebraic closures, and
universal-homeomorphism invariance for the two purely inseparable
extensions.}

\subsection{Stalkwise projective-line base change}

\begin{lemma}[Exchange on a geometric stalk]
Let $p:\mathbf P^1_S\to S$, let $g:T\to S$, and let $\bar t\to T$ have
image $\bar s\to S$.  For a torsion sheaf $F$ and every $q$, the stalk of
the specified exchange map is the canonical restriction
\[
 H^q(\mathbf P^1_{\kappa(\bar s)},F|)
 \longrightarrow
 H^q(\mathbf P^1_{\kappa(\bar t)},F|),
\]
where both residue fields are separably closed.
\end{lemma}
\deps{Definition 1.6, Proposition 7.7, and the strict-local comparison of
Proposition 6.23.}

\begin{lemma}[Conservativity of geometric stalks]
A morphism of etale sheaves of abelian groups is an isomorphism if and only
if it is an isomorphism at every geometric stalk.
\end{lemma}
\deps{Vanishing of a sheaf with all geometric stalks zero and the exactness
of stalk functors.}

\begin{proposition}[Projective-line proper base change]
For every scheme $S$, the projection
$\mathbf P^1_S\to S$ satisfies $\mathrm{PBC}$.
\end{proposition}
\deps{Lemma 7.22, Proposition 7.21, Lemma 7.23, and Proposition 2.11.}

\section{S6: finite-torsion proper base change}

\begin{theorem}[Prime-coefficient proper base change]
For every prime $r$, every proper $f:X\to Y$, every cartesian base change,
and every sheaf $F$ of $\F_r$-modules on $X_{\Et}$, the specified map
\[
 g^{-1}Rf_*F\xrightarrow{\sim}Rf'_*g'^{-1}F
\]
is an isomorphism.
\end{theorem}
\deps{Proposition 5.9 and Proposition 7.24.}

\begin{proposition}[Assembly for arbitrary torsion sheaves]
For every proper $f:X\to Y$, every cartesian base change, and every torsion
abelian sheaf $F$ on $X_{\Et}$, the exchange map of Definition 1.6 is an
isomorphism in every cohomological degree.
\end{proposition}
\deps{Theorem 8.1 and Lemmas 2.6--2.9.}

\begin{theorem}[Finite-torsion proper base change]
Let
\[
\begin{CD}
 X' @>{g'}>> X\\
 @V{f'}VV @VV{f}V\\
 Y' @>{g}>> Y
\end{CD}
\]
be cartesian with $f$ proper.  For every torsion sheaf $F$ on $X_{\Et}$,
\[
 \Ex_{\square,F}:g^{-1}Rf_*F
 \xrightarrow{\sim}Rf'_*g'^{-1}F
\]
is an isomorphism in $\Dplus(Y'_{\Et})$.  It is natural in $F$, is the
identity for identity squares, and satisfies horizontal and vertical
pasting.
\end{theorem}
\deps{Proposition 8.2 and the coherence equations of Lemma 1.8.}
\sources{Stacks Project, Tag 095T.}

\section{S7: bounded-below torsion complexes}

\begin{definition}[Hypercohomology spectral sequences]
Let $K\in\Dplus(X_{\Et})$ have torsion cohomology and choose $a$ with
$\mathcal H^q(K)=0$ for $q<a$.  Canonical truncations define
\[
 E_2^{p,q}=R^pf_*\mathcal H^q(K)
 \Longrightarrow\mathcal H^{p+q}(Rf_*K)
\]
and
\[
 E_2'{}^{p,q}=R^pf'_*g'^{-1}\mathcal H^q(K)
 \Longrightarrow\mathcal H^{p+q}(Rf'_*g'^{-1}K).
\]
The filtration in degree $n$ is the image filtration obtained from
$\tau_{\leq q}K\to K$.
\end{definition}
\deps{The standard $t$-structure, canonical truncation triangles, and a
right-derived functor.}

\begin{lemma}[Half-plane convergence]
Both spectral sequences of Definition 9.1 lie in $p\geq0$, $q\geq a$.  In
total degree $n$, only $a\leq q\leq n$ occurs.  The induced filtrations are
finite, exhaustive, and separated, so the spectral sequences converge
without a derived inverse-limit condition.
\end{lemma}
\deps{Definition 9.1 and nonnegativity of the cohomological degrees of a
right-derived functor.}

\begin{definition}[Truncation exact couples]
For $q\geq a$, apply $g^{-1}Rf_*$ and $Rf'_*g'^{-1}$ to the truncation
triangle
\[
 \tau_{\leq q-1}K\to\tau_{\leq q}K
 \to\mathcal H^q(K)[-q]\to.
\]
The exchange maps on its three terms give a morphism between the resulting
exact couples.  Its first graded component in bidegree $(p,q)$ is
\[
 g^{-1}R^pf_*\mathcal H^q(K)\longrightarrow
 R^pf'_*g'^{-1}\mathcal H^q(K).
\]
\end{definition}
\deps{Definition 1.6, Definition 9.1, and functoriality of canonical
truncation triangles.}

\begin{lemma}[Compatibility of the exact-couple differential]
Using the standard cohomological cone convention, the morphism of
Definition 9.3 commutes with the connecting maps and with every differential
of the two spectral sequences.  Its $E_2$ map is exactly the finite-torsion
exchange map for $\mathcal H^q(K)$.
\end{lemma}
\deps{Definition 9.3, naturality of cones, and Theorem 8.3.}

\begin{lemma}[Abutment identification]
In total degree $n$, the map on the abutments of the spectral sequences is
\[
 \mathcal H^n(\Ex_{\square,K}):
 g^{-1}\mathcal H^n(Rf_*K)\longrightarrow
 \mathcal H^n(Rf'_*g'^{-1}K),
\]
with the image filtrations of Definition 9.1.
\end{lemma}
\deps{Definitions 9.1 and 9.3, Lemma 9.2, and naturality of exchange in
$\tau_{\leq q}K\to K$.}

\begin{theorem}[Bounded-below derived proper base change]
For a proper cartesian square and every
$K\in\Dplus(X_{\Et})$ with torsion cohomology,
\[
 \Ex_{\square,K}:g^{-1}Rf_*K
 \xrightarrow{\sim}Rf'_*g'^{-1}K
\]
is an isomorphism in $\Dplus(Y'_{\Et})$.
\end{theorem}
\deps{Theorem 8.3 and Lemmas 9.2, 9.4, and 9.5; finite filtrations detect
an isomorphism from their associated graded objects.}
\sources{Stacks Project, Tag 0DDE.}

\section{S8: finite-level etale--pro-etale transport}

\begin{definition}[The comparison morphism of sites]
For a scheme $T$, the inclusion of etale $T$-schemes into weakly etale
$T$-schemes is continuous and preserves finite limits.  It defines a
geometric morphism
\[
 \nu_T:T_{\Proet}\longrightarrow T_{\Et},
 \qquad
 \nu_T^*:\Sh(T_{\Et})\rightleftarrows
 \Sh(T_{\Proet}):\nu_{T,*}.
\]
For a finite discrete ring $\Lambda$, the same adjunction acts on sheaves
of $\Lambda$-modules.
\end{definition}
\deps{The etale and pro-etale sites, continuity of the site inclusion, and
the construction of a geometric morphism from a continuous left-exact
functor.}

\begin{proposition}[Exactness and full faithfulness at finite level]
For finite $\Lambda$, $\nu_T^*$ is exact, its unit
\[
 F\longrightarrow\nu_{T,*}\nu_T^*F
\]
is an isomorphism for every etale $\Lambda$-module sheaf, and
\[
 R^q\nu_{T,*}\nu_T^*I=0\qquad(q>0)
\]
for every injective etale $\Lambda$-module sheaf $I$.
\end{proposition}
\deps{Definition 10.1, weakly contractible pro-etale covers, and the
acyclic-resolution criterion.}
\sources{Stacks Project, Pro-etale Cohomology, Lemma 23.1.}

\begin{theorem}[Bounded-below site comparison]
For finite $\Lambda$, derived inverse image is fully faithful on
bounded-below complexes:
\[
 \nu_T^*:\Dplus(T_{\Et},\Lambda)
 \longrightarrow\Dplus(T_{\Proet},\Lambda),
 \qquad R\nu_{T,*}\nu_T^*\xrightarrow{\sim}\operatorname{id}.
\]
\end{theorem}
\deps{Proposition 10.2 and a bounded-below injective resolution.}

\begin{definition}[Explicit etale-origin witness]
An object $K\in\Dplus(T_{\Proet},\Lambda)$ comes from the etale site if it
is equipped with a pair
\[
 (L,\alpha),\qquad
 L\in\Dplus(T_{\Et},\Lambda),\quad
 \alpha:\nu_T^*L\xrightarrow{\sim}K.
\]
The witness, rather than only its existence, is retained when maps are
transported.
\end{definition}
\deps{Theorem 10.3.}

\begin{lemma}[Comparison inverse image commutes with scheme pullback]
For $u:U\to V$ there is a natural isomorphism
\[
 \nu_U^*u_{\Et}^{-1}\xrightarrow{\sim}
 u_{\Proet}^{-1}\nu_V^*.
\]
It is compatible with identity maps, composites, and the cartesian
identification in Definition 1.3.
\end{lemma}
\deps{Definition 10.1 and the equality of the two pullback functors on
representable etale objects.}

\begin{proposition}[Comparison direct image]
Let $u:U\to V$ be quasi-compact and quasi-separated and let
$L\in\Dplus(U_{\Et},\Lambda)$.  There is a natural isomorphism
\[
 \nu_V^*Ru_{\Et,*}L
 \xrightarrow{\sim}
 Ru_{\Proet,*}\nu_U^*L.
\]
It is compatible with identities, composites, and morphisms of $L$.
\end{proposition}
\deps{Theorem 10.3, Lemma 10.5, and comparison of cohomology on pro-etale
test objects with the corresponding etale neighborhoods.}
\sources{Stacks Project, Pro-etale Cohomology, Lemma 23.1.}

\begin{lemma}[Transport identifies the exchange map]
For a cartesian square and $L\in\Dplus(X_{\Et},\Lambda)$, conjugating the
etale exchange map by Lemma 10.5 and Proposition 10.6 gives the pro-etale
unit--counit exchange map
\[
 g_{\Proet}^{-1}Rf_{\Proet,*}\nu_X^*L
 \longrightarrow
 Rf'_{\Proet,*}g_{\Proet}'{}^{-1}\nu_X^*L.
\]
\end{lemma}
\deps{Definition 1.6, Lemma 10.5, Proposition 10.6, and naturality of the
two comparison adjunctions.}

\begin{theorem}[Finite-level pro-etale proper base change]
Let $\Lambda$ be finite, let $f$ be proper, and let
$K\in\Dplus(X_{\Proet},\Lambda)$ carry an etale-origin witness
$(L,\alpha)$.  Then the pro-etale exchange map
\[
 g^{-1}Rf_*K\xrightarrow{\sim}Rf'_*g'^{-1}K
\]
is an isomorphism.  It is independent of replacing the witness by an
isomorphic witness.
\end{theorem}
\deps{Theorem 9.6, Definition 10.4, and Lemma 10.7.}

\section{S9: derived completion and inverse limits}

\begin{definition}[Tower category and adic coefficient datum]
For a sheaf of rings $B_T$ on $T_{\Proet}$, define the Grothendieck abelian
category
\[
 \operatorname{Tow}_{B}(T)=
 \operatorname{Fun}\bigl(\mathbf N_{\geq1}^{\mathrm{op}},
              \Mod(T_{\Proet},B_T)\bigr).
\]
Its objects are strict inverse sequences
$M_1\leftarrow M_2\leftarrow\cdots$, and kernels, cokernels, and exactness
are computed levelwise.  Write
$\operatorname{ev}_n:\operatorname{Tow}_{B}(T)\to
\Mod(T_{\Proet},B_T)$ for the exact evaluation functor.

Let $A$ be noetherian, let $I=(a_1,\dots,a_r)$, assume $A/I$ is torsion,
and put $P_n=A/I^n$.  Give $P_n$ the discrete topology and let
$(P_n)_T$ be the sheafification of
\[
 U\longmapsto\operatorname{Cont}(\pi_0(|U|),P_n)
\]
on $T_{\Proet}$.  Define
\[
 A_T=\varprojlim_n(P_n)_T,
 \qquad
 D(T_{\Proet},A)=D(\Mod(T_{\Proet},A_T)),
\]
and assume each $P_n$ is finite as a set and perfect as an $A$-module.  For
$K\in D(T_{\Proet},A)$ choose a functorial $K$-flat representative
$QK$.  The strict complexes
\[
 QK\otimes_{A_T}(P_n)_T
\]
and the quotient maps $P_{n+1}\to P_n$ define an object
$K_\bullet\in D(\operatorname{Tow}_{A}(T))$.  Its $n$-th evaluation is
canonically
\[
 K_n=K\otimes^L_{A_T}(P_n)_T,
\]
and it carries functorial normalization isomorphisms
\[
 K_{n+1}\otimes^L_{P_{n+1}}P_n\xrightarrow{\sim}K_n
\]
satisfying the quotient cocycle.
\end{definition}
\deps{Definitions 1.11--1.12, Grothendieck functor categories,
functorial $K$-flat resolutions, and the quotient tower of a finitely
generated ideal.}

\begin{definition}[Sequential derived inverse limit]
Let
\[
 \varprojlim:\operatorname{Tow}_{A}(T)
 \longrightarrow\Mod(T_{\Proet},A_T)
\]
be the left exact inverse-limit functor, right adjoint to the exact constant
tower functor.  Its total right derived functor is the ordinary
triangulated functor
\[
 \Rlim:D(\operatorname{Tow}_{A}(T))
 \longrightarrow D(T_{\Proet},A).
\]
Explicitly, choose a functorial $K$-injective resolution
$M_\bullet\to I_\bullet$ in the complex category of
$\operatorname{Tow}_{A}(T)$ and put
\[
 \Rlim M_\bullet=\varprojlim I_\bullet.
\]
Let $\operatorname{Fam}_{A}(T)$ be the functor category indexed by the
discrete set $\mathbf N_{\geq1}$.  The derived product
\[
 R\!\prod_nM_n
\]
means the total right derived functor of
$\prod:\operatorname{Fam}_{A}(T)\to\Mod(T_{\Proet},A_T)$ applied to the
underlying family of $M_\bullet$; it is represented by
$\prod_n\operatorname{ev}_n(I_\bullet)$.
\end{definition}
\deps{Definition 11.1, $K$-injective resolutions in Grothendieck abelian
categories, and total right derived functors of left exact right adjoints.}

\begin{proposition}[Roos triangle and Milnor exact sequence]
For every $M_\bullet\in D(\operatorname{Tow}_{A}(T))$, the transition maps
define an endomorphism $s$ of $R\!\prod_nM_n$, and there is a functorial
distinguished Roos triangle
\[
 \Rlim M_\bullet\longrightarrow R\!\prod_nM_n
 \xrightarrow{1-s}R\!\prod_nM_n
 \longrightarrow(\Rlim M_\bullet)[1].
\]
If $U$ is a weakly contractible quasi-compact object of $T_{\Proet}$, then
for every $i\in\mathbf Z$ this triangle induces the functorial exact
sequence
\[
 0\to\varprojlim{}^1_nH^{i-1}(U,M_n)
 \to H^i(U,\Rlim M_\bullet)
 \to\varprojlim_nH^i(U,M_n)\to0.
\]
\end{proposition}
\deps{Definitions 11.1--11.2, the Roos resolution for a countable inverse
system, and exactness and product preservation of sections on a weakly
contractible quasi-compact pro-etale object.}

\begin{lemma}[Mittag--Leffler vanishing]
A tower of finite abelian groups satisfies the Mittag--Leffler condition
and has $\varprojlim{}^1=0$.  The Milnor map of Proposition 11.3 is
therefore an isomorphism for such a tower.
\end{lemma}
\deps{Stabilization of descending chains in a finite group and the
product-cokernel description of $\varprojlim{}^1$.}

\begin{definition}[Intrinsic derived completion]
For $m\geq1$, let
\[
 Q_m=K(A_T;a_1^m,\dots,a_r^m)
   =\bigotimes_{j=1}^r[A_T\xrightarrow{a_j^m}A_T]
\]
in cohomological degrees $[-r,0]$.  For each factor define a chain map
\[
 [A_T\xrightarrow{a_j^{m+1}}A_T]
 \longrightarrow[A_T\xrightarrow{a_j^m}A_T]
\]
whose degree $-1$ component is multiplication by $a_j$ and whose degree
$0$ component is the identity.  Their tensor product is a chain map
$q_{m+1,m}:Q_{m+1}\to Q_m$; the composites
$q_{m+2,m}=q_{m+1,m}q_{m+2,m+1}$ make $(Q_m)_m$ a strict tower.

For a functorial $K$-flat representative $QK$, the complexes
$QK\otimes_{A_T}Q_m$ form an object of
$D(\operatorname{Tow}_{A}(T))$.  Define
\[
 K^{\wedge}_{\mathrm{Kos}}
   =\Rlim_m(QK\otimes_{A_T}Q_m).
\]
The compatible degree-zero maps $A_T\to Q_m$ define
$K\to K^{\wedge}_{\mathrm{Kos}}$.  The object $K$ is derived
$I$-complete if this map is an isomorphism.
\end{definition}
\deps{Definitions 11.1--11.2, tensor products of two-term Koszul
complexes, and the displayed chain-map identities.}

\begin{proposition}[Koszul and quotient reconstruction]
For the strict quotient tower $K_\bullet$ of Definition 11.1 there is a
functorial isomorphism
\[
 K^{\wedge}_{\mathrm{Kos}}\xrightarrow{\sim}\Rlim K_\bullet
\]
carrying the Koszul completion map to the quotient-tower completion map
$K\to\Rlim K_\bullet$.  Every derived $I$-complete object consequently
satisfies
\[
 K\xrightarrow{\sim}\Rlim K_\bullet.
\]
\end{proposition}
\deps{Definitions 11.1, 11.2, and 11.5, Artin--Rees for the negative
Koszul cohomology, and cofinality of the ideals
$(a_1^m,\dots,a_r^m)$ and $I^n$.}
\sources{Stacks Project, Tag 099M.}

\begin{lemma}[Finite reduction of reconstruction]
For every $a\geq1$, the derived counit from inverse limit to evaluation
induces a functorial isomorphism
\[
 (\Rlim K_\bullet)\otimes_A^LP_a\xrightarrow{\sim}K_a.
\]
Moreover, $(-)\otimes_A^LP_a$ commutes with the two derived countable
products in the Roos triangle of Proposition 11.3.
\end{lemma}
\deps{Propositions 11.3 and 11.6 and the bounded finite-projective
representative of $P_a$.}

\begin{definition}[Uniformly bounded constructible tower]
An adic tower is an object $M_\bullet\in
D(\operatorname{Tow}_{A}(T))$ equipped with normalization isomorphisms
$M_{n+1}\otimes^L_{P_{n+1}}P_n\simeq M_n$ satisfying the quotient
cocycle.  It is uniformly bounded constructible if there are integers
$d_0\leq d_1$, one finite stratification $T=\coprod_\alpha T_\alpha$, and,
on each stratum, a pro-etale covering $\{U_{\alpha j}\to T_\alpha\}_j$
with the following datum.  For each $U_{\alpha j}$ there is a bounded
complex $L_{\alpha j}$, supported in $[d_0,d_1]$, whose terms are finite
locally free $A_T|_{U_{\alpha j}}$-modules and whose reductions have finite
lisse cohomology, together with isomorphisms
\[
 M_n|_{U_{\alpha j}}\xrightarrow{\sim}
 L_{\alpha j}\otimes^L_{A_T}(P_n)_T|_{U_{\alpha j}}
\]
for every $n$.  These isomorphisms intertwine the normalization maps of
$M_\bullet$ with the quotient transition maps of the local reduction
tower.
\end{definition}
\deps{Definition 11.1, finite stratifications, pro-etale local finite
freeness, and finite-level constructibility.}

\begin{lemma}[Reconstruction of a uniform local tower]
Let $M_\bullet$ be a uniformly bounded constructible adic tower and put
$M=\Rlim M_\bullet$.  For every local model in Definition 11.8,
restriction of the derived-limit projections gives a functorial
isomorphism
\[
 L_{\alpha j}\xrightarrow{\sim}M|_{U_{\alpha j}}.
\]
Consequently $M$ is derived $I$-complete, the canonical maps
\[
 M\otimes_A^LP_n\xrightarrow{\sim}M_n
\]
are isomorphisms compatible with normalization, and $M$ has no cohomology
outside $[d_0,d_1]$.  No finiteness assertion about the groups of sections
on $U_{\alpha j}$ is required.
\end{lemma}
\deps{Definition 11.8, Propositions 11.3 and 11.6, Lemma 11.7, and descent
of isomorphisms and cohomological bounds along a pro-etale covering.}

\begin{definition}[Completed pullback]
For $g:T'\to T$, levelwise inverse image on strict tower complexes has a
total left derived functor
\[
 Lg^*_{\operatorname{Tow}}:
 D(\operatorname{Tow}_{A}(T))
 \longrightarrow D(\operatorname{Tow}_{A}(T')).
\]
Its $n$-th evaluation on the reduction tower $K_\bullet$ is
$Lg_n^*K_n$.  For a derived complete $K$, define
\[
 Lg^*_{\mathrm{comp}}K
   =\Rlim\bigl(Lg^*_{\operatorname{Tow}}K_\bullet\bigr)
 \quad\text{in }D(T'_{\Proet},A).
\]
\end{definition}
\deps{Definitions 11.1--11.2, levelwise $K$-flat resolutions, and
finite-level derived pullback.}

\begin{lemma}[Reduction of completed pullback]
The object $Lg^*_{\mathrm{comp}}K$ is derived $I$-complete, and for every
$n$ there is a functorial isomorphism
\[
 (Lg^*_{\mathrm{comp}}K)\otimes_A^LP_n
 \xrightarrow{\sim}Lg_n^*K_n,
\]
compatible with the transition maps and normalization cocycles.
\end{lemma}
\deps{Definitions 11.5 and 11.10, Proposition 11.6, and Lemma 11.7.}

\begin{definition}[Projection-formula map]
For $f:X\to Y$, $K\in D(X_{\Proet},A)$, and
$P\in D(Y_{\Proet},A)$, define
\[
 \pi_{f,K,P}:(Rf_*K)\otimes_A^LP
 \longrightarrow Rf_*(K\otimes_A^LLf^*P)
\]
as the adjoint of the tensor product of the counit
$Lf^*Rf_*K\to K$ with $Lf^*P$.
\end{definition}
\deps{The derived adjunction $Lf^*\dashv Rf_*$ and the symmetric monoidal
structure on the derived category.}

\begin{proposition}[Perfect projection formula]
For fixed $f$ and $K$, the full subcategory of objects $P$ for which
$\pi_{f,K,P}$ is an isomorphism is a thick triangulated subcategory
containing every finite free $A$-module.  In particular, it contains every
perfect complex, and
\[
 (Rf_*K)\otimes_A^LP_n\xrightarrow{\sim}Rf_*K_n.
\]
\end{proposition}
\deps{Definition 11.12, additivity and exactness of both functors, and the
perfectness hypothesis in Definition 11.1.}

\begin{proposition}[Direct image, derived limits, and completeness]
Levelwise direct image on tower categories has a total right derived
functor
\[
 Rf_*^{\operatorname{Tow}}:
 D(\operatorname{Tow}_{A}(X))
 \longrightarrow D(\operatorname{Tow}_{A}(Y)),
\]
with
$\operatorname{ev}_nRf_*^{\operatorname{Tow}}M_\bullet
\simeq Rf_*M_n$.  For every $M_\bullet$ the canonical map
\[
 Rf_*\Rlim M_\bullet
 \xrightarrow{\sim}
 \Rlim\bigl(Rf_*^{\operatorname{Tow}}M_\bullet\bigr)
\]
is an isomorphism.  Under the two Roos triangles, it is induced by the
canonical isomorphisms
\[
 Rf_*R\!\prod_nM_n\xrightarrow{\sim}
 R\!\prod_nRf_*M_n
\]
and their commutation with $1-s$.
If $K\in D(X_{\Proet},A)$ is derived $I$-complete,
then $Rf_*K$ is derived $I$-complete.  For every $n\geq1$, the
projection-formula morphism is a functorial isomorphism
\[
 (Rf_*K)\otimes_A^LP_n
 \xrightarrow{\sim}Rf_*(K\otimes_A^LP_n)=Rf_*K_n,
\]
and these isomorphisms commute with $P_{n+1}\to P_n$.  Under them, the
completion map for $Rf_*K$ is the composite
\[
 Rf_*K\xrightarrow{\sim}Rf_*\Rlim K_\bullet
 \xrightarrow{\sim}\Rlim Rf_*^{\operatorname{Tow}}K_\bullet.
\]
\end{proposition}
\deps{Definitions 11.1--11.2, Proposition 11.3, Proposition 11.6, the
perfect projection formula of Proposition 11.13, and the derived
adjunction $Lf^*\dashv Rf_*$.}

\begin{lemma}[Levelwise isomorphisms and complete cones]
A morphism in $D(\operatorname{Tow}_{A}(T))$ whose evaluation at every
$n$ is an isomorphism is itself an isomorphism and remains an isomorphism
after $\Rlim$.  If a derived complete object $C$ satisfies
$C\otimes_A^LP_n=0$ for every $n$, then $C=0$.  Consequently a morphism
between derived complete objects is an isomorphism once every finite
reduction is an isomorphism.
\end{lemma}
\deps{Definitions 11.1--11.2, conservativity of the evaluation functors,
Proposition 11.6, and the distinguished triangle of the cone of a
morphism.}

\begin{lemma}[Perfectness of the standard quotients]
For $A=\Zl$ and $I=(\ell)$, the quotient
$P_n=\mathbf Z/\ell^n\mathbf Z$ is perfect over $\Zl$, represented by
\[
 [\Zl\xrightarrow{\ell^n}\Zl]
\]
in degrees $-1$ and $0$.  The analogous two-term resolution holds for
$\mathcal O_E/\varpi^n$ over the ring of integers $\mathcal O_E$ of a
finite extension $E/\Ql$.
\end{lemma}
\deps{A uniformizer is a non-zero-divisor in a complete discrete valuation
ring.}

\section{S10: complete, rational, and algebraically closed coefficients}

\subsection{The complete pro-etale rectangle}

\begin{definition}[Four finite-level corners]
Let $\square$ be cartesian with $f$ proper, let $K$ be derived
$I$-complete, put $K_n=K\otimes_A^LP_n$, and put
\[
 E=Lg'^*K,
 \qquad E_n=Lg_n'^*K_n
       \xleftarrow{\sim}E\otimes_A^LP_n.
\]
Define four objects of $D(Y'_{\Proet},P_n)$:
\[
\begin{aligned}
 A_n&=Lg_n^*((Rf_*K)\otimes_A^LP_n),&
 B_n&=Lg_n^*Rf_*K_n,\\
 C_n&=(Rf'_*E)\otimes_A^LP_n,&
 D_n&=Rf'_*E_n.
\end{aligned}
\]
\end{definition}
\deps{Definitions 1.3, 1.12, 11.1, and 11.10 and compatibility of ordinary
derived pullback with finite reduction.}

\begin{definition}[Four finite-level arrows]
Let
\[
 h_n:A_n\xrightarrow{\sim}B_n,
 \qquad h'_n:C_n\xrightarrow{\sim}D_n
\]
be the perfect projection-formula isomorphisms.  Let
\[
 v_n:A_n\to C_n
\]
be the tensor reduction of the ambient ringed exchange map
\[
 \Ex^A_{\square,K}:Lg^*Rf_*K\longrightarrow Rf'_*Lg'^*K,
\]
transported through the canonical isomorphism
$Lg_n^*((Rf_*K)\otimes_A^LP_n)
\simeq(Lg^*Rf_*K)\otimes_A^LP_n$.  Let
\[
 v'_n:B_n\xrightarrow{\sim}D_n
\]
be the finite-level pro-etale exchange isomorphism of Theorem 10.8.
\end{definition}
\deps{Definitions 1.13 and 12.1, Proposition 11.13, and Theorem 10.8.}

\begin{lemma}[Finite-level exchange rectangle]
For every $n$, the diagram
\[
\begin{CD}
 A_n @>{h_n}>> B_n\\
 @V{v_n}VV @VV{v'_n}V\\
 C_n @>{h'_n}>> D_n
\end{CD}
\]
commutes.  Equivalently, $v'_nh_n=h'_nv_n$.
\end{lemma}
\deps{Definition 12.2 and naturality of the projection-formula map with
respect to the base-change mate.}

\begin{lemma}[Compatibility with the coefficient tower]
The transition morphisms $P_{n+1}\to P_n$ make each corner of Lemma 12.3
an object
\[
 A_\bullet,B_\bullet,C_\bullet,D_\bullet
 \in D(\operatorname{Tow}_{A}(Y'))
\]
after restriction of scalars, and make $h,h',v,v'$ morphisms in this
derived tower category.  Evaluation at $n$ recovers the objects and maps
of Definitions 12.1--12.2.  The rectangle equation commutes with every
transition and with every compatible endomorphism of $K$.
\end{lemma}
\deps{Lemma 12.3, naturality of exchange, and the coefficient compatibility
of the perfect projection formula.}

\begin{lemma}[The rectangle passes to the derived limit]
Put
\[
 A_\infty=\Rlim A_\bullet,\quad
 B_\infty=\Rlim B_\bullet,\quad
 C_\infty=\Rlim C_\bullet,\quad
 D_\infty=\Rlim D_\bullet,
\]
and let $H,V,H',V'$ be the images under the triangulated functor $\Rlim$
of the four tower morphisms.  Then
\[
 V'H=H'V:A_\infty\longrightarrow D_\infty.
\]
The $n$-th component of both sides is the equality of Lemma 12.3.
\end{lemma}
\deps{Definitions 11.2 and 12.1--12.2 and Lemma 12.4.}

\begin{definition}[Completed base-change map]
Assume that every $K_n$ carries a specified etale-origin witness.  Define
\[
 \Ex^{\mathrm{comp}}_{\square,K}:
 Lg^*_{\mathrm{comp}}Rf_*K
 \longrightarrow Rf'_*Lg'^*_{\mathrm{comp}}K
\]
as the ordered composite
\[
\begin{aligned}
 Lg^*_{\mathrm{comp}}Rf_*K
 &\xrightarrow{\sim}\Rlim B_\bullet\\
 &\xrightarrow{\Rlim(v'_\bullet)}\Rlim D_\bullet\\
 &\xrightarrow{\sim}
 Rf'_*\Rlim\bigl(Lg'^*_{\operatorname{Tow}}K_\bullet\bigr)
 =Rf'_*Lg'^*_{\mathrm{comp}}K.
\end{aligned}
\]
The first isomorphism is reconstruction followed by the perfect projection
formula, and the last is the inverse of direct-image/limit commutation.
\end{definition}
\deps{Theorem 10.8, Propositions 11.6, 11.13, and 11.14, and Lemma 12.5.}
\sources{Stacks Project, Pro-etale Cohomology, Theorem 30.1, Tag 09C9,
exchange-map construction.}

\begin{lemma}[Finite reduction of the completed map]
For every $n$, reduction of $\Ex^{\mathrm{comp}}_{\square,K}$ modulo
$I^n$, conjugated by the canonical reduction isomorphisms, is exactly
\[
 \Ex_{\square,K_n}:Lg_n^*Rf_*K_n
 \xrightarrow{\sim}Rf'_*Lg_n'^*K_n.
\]
These reduction squares commute with $n+1\to n$ and with compatible
endomorphisms.
\end{lemma}
\deps{Definition 12.6, Lemmas 11.7 and 11.11, and the component equation
of Lemma 12.5.}

\begin{lemma}[Coherence of completed exchange]
Completed exchange is natural in $K$, equals the identity for identity
squares, and satisfies horizontal and vertical pasting after the canonical
associativity isomorphisms.
\end{lemma}
\deps{Lemma 12.7, finite-level coherence from Lemma 1.8, and preservation
of compatible equalities by $\Rlim$.}

\begin{theorem}[Complete pro-etale proper base change]
Let $A$ and $I$ satisfy Definition 11.1, let $f:X\to Y$ be proper, and let
$g:Y'\to Y$ be arbitrary.  Suppose $K\in D(X_{\Proet},A)$ is derived
$I$-complete and, for every $n$, $K_n$ is equipped with an etale-origin
witness
\[
 L_n\in\Dplus(X_{\Et},P_n),
 \qquad \nu_X^*L_n\xrightarrow{\sim}K_n,
\]
compatible with transition maps.  Then
\[
 \Ex^{\mathrm{comp}}_{\square,K}:
 Lg^*_{\mathrm{comp}}Rf_*K
 \xrightarrow{\sim}Rf'_*Lg'^*_{\mathrm{comp}}K
\]
is an isomorphism, natural in $K$ and coherent for identities and pasted
squares.
\end{theorem}
\deps{Definition 12.6, Theorem 10.8, Proposition 11.14, Lemma 11.15, and
Lemma 12.8.}
\sources{Stacks Project, Pro-etale Cohomology, Theorem 30.1, Tag 09C9.}

\begin{definition}[Bounded constructible complete category]
Assume $T$ is noetherian of finite dimension and the torsion orders are
invertible on $T$.  Define $\Dbc(T_{\Proet},A)$ to be the full subcategory
of derived $I$-complete objects whose reduction tower has compatible etale
origins and is uniformly bounded constructible in the sense of Definition
11.8.
\end{definition}
\deps{Definitions 10.4, 11.5, and 11.8.}

\begin{proposition}[Preservation before constructible restriction]
Let $\square$ be the cartesian square of Definition 1.3.  Assume that
$X,Y,X',Y'$ are noetherian of finite dimension, that $f$ and $f'$ are
proper, and that every integer annihilating $A/I$ is invertible on all four
schemes.  Assume moreover that there is an integer $c\geq0$ such that, for
every $n\geq1$, finite-level direct image and pullback are exact
triangulated functors
\[
\begin{aligned}
 Rf_* &:D^b_c(X_{\Proet},P_n)\longrightarrow
        D^b_c(Y_{\Proet},P_n),\\
 Rf'_* &:D^b_c(X'_{\Proet},P_n)\longrightarrow
         D^b_c(Y'_{\Proet},P_n),\\
 Lg_n^* &:D^b_c(Y_{\Proet},P_n)\longrightarrow
          D^b_c(Y'_{\Proet},P_n),\\
 Lg_n'^* &:D^b_c(X_{\Proet},P_n)\longrightarrow
           D^b_c(X'_{\Proet},P_n).
\end{aligned}
\]
For $a\leq b$, write $D^b_c(Z_{\Proet},P_n)^{[a,b]}$ for the full class of
objects whose cohomology vanishes outside $[a,b]$; no triangulated structure
is asserted on this fixed-amplitude class.  Assume the containments
\[
\begin{aligned}
 Rf_*\bigl(D^b_c(X_{\Proet},P_n)^{[a,b]}\bigr)
 &\subset D^b_c(Y_{\Proet},P_n)^{[a,b+c]},\\
 Rf'_*\bigl(D^b_c(X'_{\Proet},P_n)^{[a,b]}\bigr)
 &\subset D^b_c(Y'_{\Proet},P_n)^{[a,b+c]},\\
 Lg_n^*\bigl(D^b_c(Y_{\Proet},P_n)^{[a,b]}\bigr)
 &\subset D^b_c(Y'_{\Proet},P_n)^{[a,b]},\\
 Lg_n'^*\bigl(D^b_c(X_{\Proet},P_n)^{[a,b]}\bigr)
 &\subset D^b_c(X'_{\Proet},P_n)^{[a,b]}
\end{aligned}
\]
for all $a\leq b$.  Require the four finite-level functors to commute with
the normalization isomorphisms for $P_{n+1}\to P_n$ and to carry compatible
etale-origin witnesses to compatible etale-origin witnesses.  Finally,
assume that their levelwise tower functors carry every uniformly bounded
constructible tower of Definition 11.8 to a tower of the same kind, with
local finite-free models compatible with the displayed amplitude
containments.

Then, for every $K\in\Dbc(X_{\Proet},A)$ whose reductions have a common
cohomological interval $[a,b]$, one has
\[
\begin{aligned}
 Rf_*K&\in\Dbc(Y_{\Proet},A),&
 Lg^*_{\mathrm{comp}}Rf_*K&\in\Dbc(Y'_{\Proet},A),\\
 Lg'^*_{\mathrm{comp}}K&\in\Dbc(X'_{\Proet},A),&
 Rf'_*Lg'^*_{\mathrm{comp}}K&\in\Dbc(Y'_{\Proet},A).
\end{aligned}
\]
The two pullbacks have reduction amplitude $[a,b]$ and the two direct
images have reduction amplitude contained in $[a,b+c]$.
\end{proposition}
\deps{Definition 12.10, Lemmas 11.9 and 11.11, Proposition 11.14, and the
stated finite-level preservation hypotheses.}

\begin{corollary}[Constructible complete proper base change]
Under the hypotheses of Proposition 12.11, Theorem 12.9 restricts to a
natural isomorphism of ordinary triangulated functors
\[
 Lg^*_{\mathrm{comp}}Rf_*K
 \xrightarrow{\sim}Rf'_*Lg'^*_{\mathrm{comp}}K
 \quad\text{in }\Dbc(Y'_{\Proet},A)
\]
for every $K\in\Dbc(X_{\Proet},A)$.
\end{corollary}
\deps{Theorem 12.9 and Proposition 12.11.}

\subsection{Finite coefficient fields}

\begin{definition}[Rational category from an integral category]
Let $E/\Ql$ be finite, let $\mathcal O_E$ be its ring of integers, and let
$\varpi$ be a uniformizer.  Define
\[
 \Dbc(T,E)=
 \bigl(\Dbc(T,\mathcal O_E)[1/\varpi]\bigr)^\natural.
\]
The localization has the same objects and morphisms
\[
 \Hom(K,L)[1/\varpi]
 =\Hom_{\Dbc(T,\mathcal O_E)}(K,L)\otimes_{\mathcal O_E}E.
\]
The superscript $\natural$ denotes the Karoubi envelope: an object is
$(K,e)$ with $e^2=e$, and a morphism
$u:(K,e)\to(L,d)$ satisfies $u=d\,u\,e$.
\end{definition}
\deps{Definition 12.10, localization at a central natural endomorphism,
and the Karoubi completion of a triangulated category.}

\begin{proposition}[Functorial transport to a finite field]
Every exact $\mathcal O_E$-linear functor among the integral constructible
categories, and every natural transformation between such functors,
extends uniquely to the localization and then to the Karoubi envelope.  In
particular $Rf_*$, completed pullback, and their exchange transformation
define exact functors and a natural transformation on $\Dbc(-,E)$.
\end{proposition}
\deps{Definition 12.13 and functoriality of localization and Karoubi
completion.}

\begin{theorem}[Proper base change over a finite extension]
For every finite $E/\Ql$, every proper cartesian square, and every
$K\in\Dbc(X,E)$,
\[
 Lg^*Rf_*K\xrightarrow{\sim}Rf'_*Lg'^*K
 \quad\text{in }\Dbc(Y',E).
\]
This isomorphism is independent of an integral lattice and is natural and
coherent for identities and pasted squares.
\end{theorem}
\deps{Corollary 12.12 and Proposition 12.14; an isomorphism remains an
isomorphism under localization and passage to a direct summand.}

\begin{proposition}[Compatibility with finite scalar extension]
For a finite extension $E\subset E'$, derived scalar extension is an exact
functor
\[
 T_{E'/E}=E'\otimes_E^L-:
 \Dbc(T,E)\longrightarrow\Dbc(T,E').
\]
Choose $T_{E/E}=\operatorname{id}$ literally.  For
$E\subset E'\subset E''$, derived tensor associativity gives a specified
exact natural isomorphism
\[
 a_{E'',E',E}:T_{E''/E'}T_{E'/E}
 \xRightarrow{\sim}T_{E''/E}.
\]
These associators satisfy the unit equations and the pentagon for four
nested fields.  Exactness supplies natural shift comparisons
\[
 \sigma_{E'/E,K}:T_{E'/E}(K[1])
 \xrightarrow{\sim}(T_{E'/E}K)[1]
\]
which commute with the associators.

There are coherent natural isomorphisms
\[
 E'\otimes_E^LRf_{*,E}K\xrightarrow{\sim}
 Rf_{*,E'}(E'\otimes_E^LK)
\]
and analogous isomorphisms for pullback.  Under these isomorphisms, scalar
extension of the $E$-exchange map is the $E'$-exchange map.
\end{proposition}
\deps{Proposition 12.14, the integral projection formula, and associativity
of derived scalar extension.}

\subsection{The \texorpdfstring{$\Qlb$}{Qbar-l} Karoubi colimit}

\begin{definition}[Finite-stage pair category]
Let $\mathsf{FinExt}_\ell$ be the poset of finite intermediate fields
$\Ql\subset E\subset\Qlb$; it is filtered because the compositum of two
finite fields is finite.  Put
\[
 \mathcal C_E(T)=\Dbc(T,E),
 \qquad T_{F/E}=F\otimes_E^L-:\mathcal C_E(T)\to\mathcal C_F(T).
\]
Use the normalized units, associators $a_{G,F,E}$, and shift comparisons
$\sigma_{F/E}$ of Proposition 12.16.

Define an object of $\mathcal P_T$ to be a pair $(E,K)$ with
$K\in\mathcal C_E(T)$.  A morphism representative
$(M,u):(E,K)\to(F,L)$ consists of a finite field $M\supset E,F$ and a map
\[
 u:T_{M/E}K\longrightarrow T_{M/F}L
 \quad\text{in }\mathcal C_M(T).
\]
For $M\subset N$, transport it to
\[
 u_N=
 a_{N,M,F,L}\circ T_{N/M}(u)\circ a_{N,M,E,K}^{-1}.
\]
Representatives $(M,u)$ and $(N,v)$ are equivalent if $u_P=v_P$ in
$\mathcal C_P(T)$ for some finite $P\supset M,N$.  Define
\[
 \Hom_{\mathcal P_T}((E,K),(F,L))
 =\colim_{M\supset E,F}
 \Hom_{\mathcal C_M(T)}(T_{M/E}K,T_{M/F}L)
\]
as the resulting equivalence classes.  Addition, identities, negatives,
and composition are computed after transporting all representatives to one
common finite overfield.
\end{definition}
\deps{Proposition 12.16, finiteness of composita, and filtered colimits of
abelian groups.}

\begin{proposition}[Common-stage additive and shift calculus]
The equivalence relation of Definition 12.17 is reflexive, symmetric, and
transitive.  Transport is additive and obeys $(u_N)_P=u_P$ for
$M\subset N\subset P$.  The common-stage operations are
representative-independent and make $\mathcal P_T$ an additive category.
Every finite diagram of objects, morphisms, composites, and asserted
equalities in $\mathcal P_T$ is represented over one finite field, and an
asserted equality holds exactly when it holds after passage to a further
common stage.

Define
\[
 (E,K)[1]=(E,K[1])
\]
and, for $f=[M,u]:(E,K)\to(F,L)$, define
\[
 f[1]=\left[M,
 \sigma_{M/F,L}^{-1}\circ u[1]\circ\sigma_{M/E,K}\right].
\]
This gives a well-defined additive autoequivalence $[1]$ with inverse
$[-1]$, and every finite-stage inclusion
$j_E:\mathcal C_E(T)\to\mathcal P_T$ commutes with shifts through the
specified comparison.
\end{proposition}
\deps{Definition 12.17, the unit and pentagon equations and the shift
coherence of Proposition 12.16, and the additive category laws at a common
finite stage.}

\begin{definition}[Common-stage cones and distinguished triangles]
Let $f=[M,u]:(E,K)\to(F,L)$.  Choose in $\mathcal C_M(T)$ a distinguished
cone triangle
\[
 T_{M/E}K\xrightarrow{u}T_{M/F}L
 \xrightarrow{i_u}C(u)\xrightarrow{p_u}(T_{M/E}K)[1].
\]
Its image in $\mathcal P_T$ is the triangle with cone object $(M,C(u))$,
second arrow represented by $i_u$, and third arrow represented by
\[
 \sigma_{M/E,K}^{-1}\circ p_u:
 C(u)\longrightarrow T_{M/E}(K[1]).
\]
A triangle in $\mathcal P_T$ is distinguished if it is isomorphic to such
a common-stage cone triangle.
\end{definition}
\deps{Proposition 12.18 and the distinguished cone triangles in each
ordinary triangulated category $\mathcal C_M(T)$.}

\begin{proposition}[Compatibility and independence of common-stage cones]
If $M\subset N$, exactness of $T_{N/M}$ identifies the scalar extension of
a cone triangle for $u$ with a cone triangle for $u_N$.  If two
representatives define the same morphism of $\mathcal P_T$, their cone
triangles become isomorphic after passage to a field where the
representatives agree.  Any two cone triangles selected for one
finite-stage morphism are isomorphic by a morphism of triangles whose first
two components are identities.  Thus distinguishedness in Definition
12.19 is independent of every representative, common field, and cone
choice.
\end{proposition}
\deps{Definitions 12.17 and 12.19, Proposition 12.18, exactness of scalar
extension, and the finite-stage morphism-of-triangles axiom.}

\begin{theorem}[Verdier axioms and stagewise triangulation]
The shift and distinguished triangles of Definitions 12.18--12.19 satisfy
the four Verdier axioms:
\begin{enumerate}
\item every morphism has a distinguished cone triangle, the identity
triangle is distinguished, and distinguishedness is closed under
isomorphism;
\item a triangle is distinguished if and only if its signed rotation is
distinguished;
\item every commutative square between the first arrows of two
distinguished triangles extends to a morphism of triangles;
\item cone triangles for composable maps can be placed in an octahedral
diagram whose required faces commute and whose remaining cone triangle is
distinguished.
\end{enumerate}
Consequently $\mathcal P_T$ is an ordinary triangulated category.  Every
finite-stage inclusion $j_E:\mathcal C_E(T)\to\mathcal P_T$ is exact, and
the identifications $j_FT_{F/E}\simeq j_E$ are exact natural
isomorphisms.
\end{theorem}
\deps{Propositions 12.18 and 12.20, the finite common-stage principle, and
the axioms TR1--TR4 in every $\mathcal C_M(T)$.}

\begin{definition}[$\Qlb$ coefficient category]
For $\alpha\in\Qlb$ and $f=[M,u]$ in $\mathcal P_T$, choose a finite field
$N\supset M,\Ql(\alpha)$ and define
\[
 \alpha f=[N,\alpha u_N].
\]
Define
\[
 \Dbc(T,\Qlb)=\operatorname{Kar}(\mathcal P_T).
\]
Thus an object is $((E,K),e)$ with $e^2=e$ in $\mathcal P_T$, and a
morphism $u:((E,K),e)\to((F,L),d)$ satisfies $u=d\,u\,e$; its identity is
$e$.  Define the shift by $((E,K),e)[1]=((E,K[1]),e[1])$.  A triangle is
distinguished if it is isomorphic to a direct summand of the image of a
distinguished triangle of $\mathcal P_T$.
\end{definition}
\deps{Theorem 12.21, finite-stage scalar multiplication, and the ordinary
Karoubi envelope.}

\begin{theorem}[Idempotent-complete triangulation]
The scalar operation of Definition 12.22 is independent of $N$ and of the
representative of $f$ and makes $\mathcal P_T$ $\Qlb$-linear.  The shift
and triangles of Definition 12.22 give $\Dbc(T,\Qlb)$ its canonical
idempotent-complete ordinary triangulated structure.  The canonical
functor $\mathcal P_T\to\Dbc(T,\Qlb)$ is exact and $\Qlb$-linear, and for
every finite $E$ the composite
\[
 \Dbc(T,E)\xrightarrow{j_E}\mathcal P_T
 \longrightarrow\Dbc(T,\Qlb)
\]
is exact and $E$-linear.
\end{theorem}
\deps{Definitions 12.17 and 12.22, Theorem 12.21, common-stage scalar
calculus, and the triangulated Karoubi-envelope theorem.}
\sources{Balmer--Schlichting, \emph{Idempotent completion of triangulated
categories}, Journal of Algebra 236 (2001), Theorem 1.5.}

\begin{proposition}[Descent of coefficient-compatible functors]
Let $\Phi_E:\mathcal C_E(T)\to\mathcal C_E(S)$ be exact functors and let
\[
 \beta_{F/E}:T^S_{F/E}\Phi_E
 \xRightarrow{\sim}\Phi_FT^T_{F/E}
\]
be exact natural isomorphisms satisfying the unit, associator, and shift
equations.  Then they induce an exact functor
$\Phi:\mathcal P_T\to\mathcal P_S$.  On a representative
$u:T^T_{M/E}K\to T^T_{M/F}L$, its value is represented by
\[
 T^S_{M/E}\Phi_EK\xrightarrow{\beta_{M/E}}
 \Phi_MT^T_{M/E}K\xrightarrow{\Phi_M(u)}
 \Phi_MT^T_{M/F}L\xrightarrow{\beta_{M/F}^{-1}}
 T^S_{M/F}\Phi_FL.
\]
This functor extends to the Karoubi envelopes.  Compatible natural
transformations and their equations descend by the same formula.

In particular, the scalar-compatible families $(Rf_{*,E})_E$ and
$(Lg_E^*)_E$ descend to exact functors on $\Qlb$ coefficient categories,
and the finite-stage exchange maps descend to a natural transformation.
\end{proposition}
\deps{Definitions 12.17 and 12.22, Theorems 12.21 and 12.23, and the
coherence statements of Proposition 12.16.}

\begin{theorem}[$\Qlb$ proper base change]
Let $f:X\to Y$ be proper, let $g:Y'\to Y$ be arbitrary, and let
$K\in\Dbc(X,\Qlb)$.  The descended exchange transformation is an
isomorphism
\[
 Lg^*Rf_*K\xrightarrow{\sim}Rf'_*Lg'^*K
 \quad\text{in }\Dbc(Y',\Qlb).
\]
If $K$ is represented by $((E,K_E),e)$, this map is represented after a
finite overfield $M$ by the $M$-coefficient exchange isomorphism restricted
to the direct summand defined by the transported idempotent.  The result is
independent of $E$, $M$, and the representative of the idempotent.
\end{theorem}
\deps{Theorem 12.15, Theorem 12.23, and Proposition 12.24.}

\section{Final geometric-fiber form and references}

\begin{corollary}[Finite-coefficient geometric-fiber formula]
Let $f:X\to Y$ be proper, let $\bar y:\Spec\Omega\to Y$ be geometric,
and let $i_{\bar y}:X_{\bar y}\to X$ and
$a_{\bar y}:X_{\bar y}\to\Spec\Omega$ be the projections.  Let $\Lambda$
be a finite ring and let $K\in\Dplus(X_{\Et},\Lambda)$.  The etale
base-change morphism is an isomorphism
\[
 \bar y^{-1}Rf_*K\xrightarrow{\sim}
 Ra_{\bar y,*}Li_{\bar y}^*K.
\]
For a finite $\Lambda$-module sheaf $F$ this gives, for every $q\geq0$,
\[
 (R^qf_*F)_{\bar y}\xrightarrow{\sim}
 H^q(X_{\bar y},F|_{X_{\bar y}}).
\]
\end{corollary}
\deps{Theorem 9.6 specialized to the cartesian square defined by
$\bar y$.}

\begin{corollary}[Complete-adic geometric-fiber formula]
In the setup of Corollary 13.1, let $(A,I)$ satisfy Definition 11.1 and let
$K\in D(X_{\Proet},A)$ satisfy the completeness and compatible
etale-origin hypotheses of Theorem 12.9.  Then the completed-pullback
exchange morphism is an isomorphism
\[
 L\bar y^*_{\mathrm{comp}}Rf_*K
 \xrightarrow{\sim}
 Ra_{\bar y,*}L(i_{\bar y})^*_{\mathrm{comp}}K
 \quad\text{in }D((\Spec\Omega)_{\Proet},A).
\]
Its reduction modulo $I^n$ is the finite-level pro-etale geometric-fiber
isomorphism for $K_n$.
\end{corollary}
\deps{Theorem 12.9 and Lemma 12.7, specialized to the geometric-point
base change.}

\begin{corollary}[Finite-extension geometric-fiber formula]
In the setup of Corollary 13.1, let $E/\Ql$ be finite and let
$K\in\Dbc(X,E)$.  The $E$-coefficient exchange morphism is an isomorphism
\[
 L\bar y^*Rf_*K\xrightarrow{\sim}
 Ra_{\bar y,*}Li_{\bar y}^*K
 \quad\text{in }\Dbc(\Spec\Omega,E).
\]
\end{corollary}
\deps{Theorem 12.15 specialized to the geometric-point base change.}

\begin{corollary}[$\Qlb$ geometric-fiber formula]
In the setup of Corollary 13.1, let $K\in\Dbc(X,\Qlb)$.  The descended
exchange morphism is an isomorphism
\[
 L\bar y^*Rf_*K\xrightarrow{\sim}
 Ra_{\bar y,*}Li_{\bar y}^*K
 \quad\text{in }\Dbc(\Spec\Omega,\Qlb).
\]
For a representative $((E,K_E),e)$, this is represented over a finite
overfield by the direct summand of the finite-extension formula selected
by the transported idempotent.
\end{corollary}
\deps{Theorem 12.25 specialized to the geometric-point base change.}

\paragraph{Source locators.}
The finite torsion spine is located in the Stacks Project, Etale
Cohomology, Sections 59.82--59.85 and 59.91.  Its endpoint is Tag 095T;
the bounded-below upgrade is Tag 0DDE.  The complete pro-etale endpoint is
the Stacks Project, Pro-etale Cohomology, Theorem 30.1, Tag 09C9.  The
intermediate locators used above are Tags 03QP, 0A4E, 0A4D, 0A4C, 0A4F,
03Q8, 0202, 09ZG, 0A0B, 09ZF, 09ZD, 09ZE, 09ZI, 0A51, 0A48, 0A5G,
0A5H, 03Q6, 03Q9, 03SI, 0A5C, 0A3R, 0GJ1, 03SG, 0A5B, 0A3Q, 03SC,
and 0A5E.

\end{document}
